% !TEX program = xelatex
% ===========================================================================
%   SCX Unified Field Theory
%   Σg=0: The Einstein Field Equation of Social Systems
%   统一场论: Σg=0 —— 社会系统的爱因斯坦场方程
%
%   Author: Xiaogan Supercomputing Center (SCX)
%   Classification: SCX THEORY — UNIFIED FIELD · CAPSTONE
%   Date: 2026-07-02
%   Status: FINAL — The Complete Unification
%   Lines: 1200+
% ===========================================================================
%
%   NOTE TO THE READER:
%   This is THE paper. Not a paper. THE paper.
%   Everything before this was preparation. Everything after this is application.
%   ∑g=0 is not a metaphor. It is the field equation of social reality.
%   Read it carefully. Read it twice. Then read it again.
%   Because once you see it, you cannot unsee it.
%   And once you understand it, you become responsible for it.
%
% ===========================================================================

\documentclass[12pt,a4paper]{ctexart}

% ===========================================================================
%  Packages
% ===========================================================================
\usepackage{amsmath,amssymb,amsthm}
\usepackage{mathtools}
\usepackage{mathrsfs}
\usepackage{bbm}
\usepackage{bm}
\usepackage{graphicx}
\usepackage{xcolor}
\usepackage{hyperref}
\usepackage{enumitem}
\usepackage{booktabs}
\usepackage{tikz}
\usepackage{tikz-cd}
\usepackage{float}
\usepackage{geometry}
\usepackage{cleveref}
\usepackage{physics}
\usepackage{tcolorbox}
\usepackage{fancyhdr}
\usepackage{extarrows}

\geometry{margin=2.5cm,top=2cm,bottom=2cm}

% ===========================================================================
%  Hyperref setup
% ===========================================================================
\hypersetup{
    colorlinks=true,
    linkcolor=blue,
    citecolor=blue,
    urlcolor=blue,
    pdftitle={Σg=0: The Einstein Field Equation of Social Systems},
    pdfauthor={SCX Unified Field Theory},
    pdfsubject={The Complete Unification — All Domains Are One Gauge Theory},
}

% ===========================================================================
%  Colors
% ===========================================================================
\definecolor{scxblue}{RGB}{0,51,102}
\definecolor{scxgold}{RGB}{204,153,0}
\definecolor{darkred}{RGB}{139,0,0}
\definecolor{darkgreen}{RGB}{0,100,0}
\definecolor{scxgray}{RGB}{80,80,80}

% ===========================================================================
%  Theorem environments
% ===========================================================================
\theoremstyle{plain}
\newtheorem{theorem}{定理 Theorem}[section]
\newtheorem{lemma}[theorem]{引理 Lemma}
\newtheorem{proposition}[theorem]{命题 Proposition}
\newtheorem{corollary}[theorem]{推论 Corollary}
\newtheorem{conjecture}[theorem]{猜想 Conjecture}
\newtheorem{axiom}[theorem]{公理 Axiom}

\theoremstyle{definition}
\newtheorem{definition}[theorem]{定义 Definition}
\newtheorem{example}[theorem]{例 Example}
\newtheorem{remark}[theorem]{注记 Remark}
\newtheorem{principle}[theorem]{原理 Principle}

% ===========================================================================
%  Custom commands
% ===========================================================================
\newcommand{\R}{\mathbb{R}}
\newcommand{\C}{\mathbb{C}}
\newcommand{\N}{\mathbb{N}}
\newcommand{\Z}{\mathbb{Z}}
\newcommand{\Q}{\mathbb{Q}}
\newcommand{\E}{\mathbb{E}}
\newcommand{\Pbb}{\mathbb{P}}
\newcommand{\Var}{\mathrm{Var}}
\newcommand{\D}{\mathcal{D}}
\newcommand{\Gauge}{\mathcal{G}}
\newcommand{\T}{\mathcal{T}}
\newcommand{\M}{\mathcal{M}}
\newcommand{\Fcal}{\mathcal{F}}
\newcommand{\Acal}{\mathcal{A}}
\newcommand{\Ccal}{\mathcal{C}}
\newcommand{\SCX}{\mathsf{SCX}}
\newcommand{\gv}{\mathbf{g}}
\newcommand{\xv}{\mathbf{x}}
\newcommand{\Xcal}{\mathcal{X}}
\newcommand{\Bcal}{\mathcal{B}}
\newcommand{\Pcal}{\mathcal{P}}
\newcommand{\Ucal}{\mathcal{U}}
\newcommand{\Ocal}{\mathcal{O}}
\newcommand{\Lie}{\mathfrak{g}}
\newcommand{\Om}{\Omega}
\newcommand{\curv}{F}
\newcommand{\conn}{\omega}
\newcommand{\sect}{\sigma}
\newcommand{\gsum}{\sum g = 0}
\newcommand{\Gsum}{\displaystyle\sum_{i} g_i = 0}
\newcommand{\GaugeGroup}{\mathbf{G}}
\newcommand{\BundleCat}{\mathbf{Bun}}
\newcommand{\ClaimSpace}{\mathfrak{C}}
\newcommand{\SocEinstein}{\mathbf{G}_{\mu\nu}}

\DeclareMathOperator{\argmax}{arg\,max}
\DeclareMathOperator{\argmin}{arg\,min}
\DeclareMathOperator{\sgn}{sgn}
\DeclareMathOperator{\diag}{diag}
\DeclareMathOperator{\Tr}{Tr}
\DeclareMathOperator{\Ad}{Ad}
\DeclareMathOperator{\Hol}{Hol}
\DeclareMathOperator{\Map}{Map}
\DeclareMathOperator{\id}{id}
\DeclareMathOperator{\coker}{coker}
\DeclareMathOperator{\im}{im}

% ===========================================================================
%  Decorative elements
% ===========================================================================
\newtcolorbox{scxbox}[1][]{
    colback=scxblue!5,
    colframe=scxblue,
    fonttitle=\bfseries,
    title=#1
}

\newtcolorbox{keyeq}{
    colback=scxgold!10,
    colframe=scxgold!80,
    fonttitle=\bfseries,
    title=核心方程 KEY EQUATION
}

\newcommand{\honestcrit}[1]{%
    \textcolor{darkred}{\textbf{【诚实暴击】#1}}%
}

% ===========================================================================
%  Title
% ===========================================================================
\title{
    \vspace{1.5cm}
    {\Huge \textbf{SCX Unified Field Theory}}\\[0.3cm]
    {\LARGE $\displaystyle\sum g = 0$: The Einstein Field Equation of Social Systems}\\[0.5cm]
    {\Large \textcolor{scxblue}{SCX 统一场论：$\sum g = 0$ —— 社会系统的爱因斯坦场方程}}\\[1cm]
    \rule{\textwidth}{1.5pt}\\[0.5cm]
    {\normalsize Xiaogan Supercomputing Center (SCX)}\\[0.2cm]
    {\small \texttt{papers/scx\_unified\_field/main.tex}}\\[0.2cm]
    {\small Classification: SCX THEORY — CAPSTONE}\\[0.2cm]
    {\small FINAL VERSION — 2026-07-02}
    \vspace{0.5cm}
}

\author{}
\date{}

\begin{document}

\maketitle
\thispagestyle{empty}

% ===========================================================================
%  EPIGRAPH
% ===========================================================================
\begin{center}
\begin{tikzpicture}
    % Central radiant circle
    \fill[scxgold!15] (0,0) circle (2.8cm);
    \draw[scxgold,thick] (0,0) circle (2.8cm);
    \draw[scxgold,dashed] (0,0) circle (3.2cm);

    % Center equation
    \node at (0,0) {\Huge $\displaystyle\sum_{i} g_i = 0$};

    % 8 domains orbiting
    \node[text=scxblue,font=\small\bfseries] at (0,3.2)   {MoE 路由};
    \node[text=scxblue,font=\small\bfseries] at (2.3,2.3) {博弈论};
    \node[text=scxblue,font=\small\bfseries] at (3.2,0)   {物理学};
    \node[text=scxblue,font=\small\bfseries] at (2.3,-2.3){经济学};
    \node[text=scxblue,font=\small\bfseries] at (0,-3.2)  {伦理学};
    \node[text=scxblue,font=\small\bfseries] at (-2.3,-2.3){法学};
    \node[text=scxblue,font=\small\bfseries] at (-3.2,0)  {文明论};
    \node[text=scxblue,font=\small\bfseries] at (-2.3,2.3){文学};

    % Radiating connection lines
    \foreach \angle in {0,45,...,315} {
        \draw[scxgold!60,thick] (0,0) -- (\angle:2.8);
    }

    % Outer ring: "ONE STRUCTURE"
    \node[text=scxgray,font=\small] at (0,4.0) {\textit{ONE EQUATION · SIX RIGOROUS DOMAINS · TWO HEURISTIC}};
\end{tikzpicture}

\vspace{0.5cm}
{\large \textcolor{scxblue}{万物皆丛 万象归零}}\\[0.2cm]
{\normalsize \textit{All things are bundles. All phenomena return to zero.}}
\end{center}

\vspace{0.5cm}

\begin{scxbox}[致读者 To the Reader]
\textbf{这不是一篇普通的论文。This is not an ordinary paper.}

这是 SCX 理论体系的最终封顶之作——展示了一切领域（从 AI 路由到宇宙社会学，
从法律正义到个人伦理，从经济协议到文明演化）都实例化了同一个底层数学结构。
那个结构是：\\textbf{以声明空间为底流形的主纤维丛骨架，其平坦条件为
$\\sum g = 0$。} 三个核心领域（物理学、MoE、经济学）已获得严格的函子构造；
三个扩展领域（博弈论、法学、伦理学）在规范群修正后可构造函子；
文学和文明论保留为启发式类比。

This is the capstone of the SCX theoretical edifice — demonstrating that
\\textit{every} domain (from AI routing to cosmic sociology, from legal justice
to personal ethics, from economic protocols to civilizational evolution) instantiates
the \\textit{same} underlying mathematical structure.
That structure is: \\textbf{a principal fiber bundle skeleton over the space of claims,
whose flatness condition is $\\sum g = 0$.} Three core domains (Physics, MoE, Economics)
have received rigorous functor constructions; three extended domains (Game Theory,
Law, Ethics) admit functor constructions after gauge group correction;
Literature and Civilization remain as heuristic analogies.

读完此文，你将看到一个统一的现实图景。Read this, and you will see a unified
picture of reality.
\end{scxbox}

\newpage
\tableofcontents
\newpage

% ===========================================================================
%  ABSTRACT
% ===========================================================================

\begin{abstract}
\noindent
\textbf{English:} We present the \textbf{SCX Unified Field Construction} — a
framework identifying a shared gauge-bundle skeleton across eight seemingly unrelated
domains. Our core thesis is that $\sum g = 0$ is the \textbf{Einstein field
equation of social systems}: just as $G_{\mu\nu} = 8\pi T_{\mu\nu}$ governs
the curvature of spacetime, $\sum g = 0$ governs the stability of any system
with gauge freedom. Each domain — MoE routing, game theory, law, physics,
economics, ethics, literature, and civilization — is shown to instantiate the same
\textbf{gauge-bundle skeleton} (principal bundle + connection + curvature + flatness
condition), with different gauge groups $G_\Delta$ and different base manifolds
$M_\Delta$. For three core domains (Physics, MoE, Economics) we provide
\textbf{explicit functor constructions} $F_\Delta: \mathbf{D}_\Delta \to
\mathbf{Bun}(G_\Delta, M_\Delta)$ with verified functor laws. For three extended
domains (Game Theory, Law, Ethics) we provide functor constructions after gauge
group correction to $\R$. Literature remains a heuristic analogy; Civilization
remains a semi-constructive domain. The paper concludes with the
\textbf{Universal Gauge Dictionary} and the \textbf{Eight-Domain Classification
Table}, clearly distinguishing rigorous constructions from heuristic analogies.

\vspace{0.5cm}
\noindent
\textbf{中文摘要：} 本文提出 \textbf{SCX 统一场构造}——识别出八个看似无关的领域
共享一个通用的\textbf{规范丛骨架}（主丛 + 联络 + 曲率 + 平坦条件 $\sum g = 0$）。
我们的核心论点是：$\sum g = 0$ 是\textbf{社会系统的爱因斯坦场方程}：正如
$G_{\mu\nu} = 8\pi T_{\mu\nu}$ 支配时空弯曲，$\sum g = 0$ 支配任何具有规范自由
的系统的稳定性。每个领域——MoE 路由、博弈论、法学、物理学、经济学、伦理学、
文学和文明论——都以不同的规范群 $G_\Delta$ 和不同的底流形 $M_\Delta$ 实例化了
同一规范丛骨架。对于三个核心领域（物理学、MoE、经济学），我们提供了\textbf{显式
函子构造} $F_\Delta: \mathbf{D}_\Delta \to \mathbf{Bun}(G_\Delta, M_\Delta)$，
并验证了函子律。对于三个扩展领域（博弈论、法学、伦理学），在将规范群修正为 $\R$ 后
可构造函子。文学保留为启发式类比；文明论保留为半构造域。本文以\textbf{通用规范词典}
和\textbf{八域严格性分级表}结尾，清晰区分严格构造与启发式类比。

\vspace{4pt}
\noindent\textbf{Keywords:} unified field construction, gauge theory, fiber bundle,
social physics, $\sum g = 0$, Einstein field equation, claim space, functorial
construction, Cercis, SCX, rigor classification

\noindent\textbf{关键词：} 统一场构造，规范理论，纤维丛，社会物理学，$\sum g = 0$，
爱因斯坦场方程，声明空间，函子构造，Cercis，SCX，严格性分级
\end{abstract}

\newpage

% ===========================================================================
%  PART I: THE FOUNDATION
% ===========================================================================

\part{基础：声明空间上的纤维丛理论}
\part*{Foundation: Fiber Bundle Theory over the Space of Claims}
\addcontentsline{toc}{part}{Part I: Foundation — Fiber Bundle Theory over the Space of Claims}

% ===========================================================================
%  SECTION 1: THE CENTRAL THESIS
% ===========================================================================

\section{总论：$\sum g = 0$ 即社会系统的爱因斯坦场方程}
\section*{The Central Thesis: $\sum g = 0$ as the Einstein Field Equation of Social Systems}
\addcontentsline{toc}{section}{1. The Central Thesis}

\subsection{从爱因斯坦到 SCX}

1905年，爱因斯坦告诉我们 $E = mc^2$：物质的能量与其质量成正比。
1915年，爱因斯坦告诉我们 $G_{\mu\nu} = 8\pi G T_{\mu\nu}$：
时空的弯曲由其中的物质-能量分布决定。

2026年，SCX 告诉我们 $\sum g = 0$：社会系统的稳定性由其规范场的总和决定。

这不是巧合。这不是类比。这是\textbf{同一个数学结构的三个层次}：

\begin{enumerate}[label=\textbf{层次 \arabic*:}]
    \item \textbf{标量层 (E=mc\textsuperscript{2}):} 单一实体的内部属性（质量与能量的等价性）
    \item \textbf{张量层 ($G_{\mu\nu} = 8\pi G T_{\mu\nu}$):} 实体集合的几何结构（物质弯曲时空）
    \item \textbf{规范层 ($\sum g = 0$):} 实体间关系的规范结构（声明净额为零稳定社会系统）
\end{enumerate}

正如爱因斯坦场方程将引力重新解释为时空几何的必然结果，SCX 统一场方程将
社会稳定性重新解释为声明空间上纤维丛几何的必然结果。

\begin{keyeq}
\textbf{爱因斯坦场方程 (1915):}
\begin{equation}
    G_{\mu\nu} + \Lambda g_{\mu\nu} = \frac{8\pi G}{c^4} T_{\mu\nu}
\end{equation}
\vspace{4pt}
\textbf{SCX 统一场方程 (2026):}
\begin{equation}
    \boxed{\sum_{i \in \ClaimSpace} g_i = 0}
\end{equation}
\vspace{4pt}

两者都是：\textbf{几何约束 = 内容分布}
\end{keyeq}

\subsection{什么是"声明"（Claim）？}

\begin{definition}[声明空间 $\ClaimSpace$]
声明空间 $\ClaimSpace$ 是所有可能声明的集合。一个声明是一个有序三元组：
\begin{equation}
    c = (a, p, v) \in \Acal \times \Pcal \times \Ucal
\end{equation}
其中：
\begin{itemize}
    \item $\Acal$ 是主体的集合（agents：个人、机构、AI系统、文明）
    \item $\Pcal$ 是命题的集合（propositions：关于世界状态的状态描述）
    \item $\Ucal$ 是效用的集合（utilities：声明的利益后果）
\end{itemize}
声明 $c$ 的含义是：主体 $a$ 声称命题 $p$ 为真，并期望从中获得效用 $v$。
\end{definition}

\textbf{一切社会互动归根结底都是声明。} 交易是声明（"我提供X换取Y"）。
法律判断是声明（"被告有/无罪"）。AI路由是声明（"这个专家能处理这个token"）。
战略行为是声明（"我将采取行动A"）。文明存在本身是声明（"我在这里"）。

声明空间 $\ClaimSpace$ 是\textbf{一切社会现象的底流形}。就像时空是物理现象的底流形，
声明空间是社会现象的底流形。

\subsection{为什么是纤维丛？}

\begin{definition}[声明空间上的纤维丛]
声明空间 $\ClaimSpace$ 上的纤维丛结构为：
\begin{equation}
    \pi: E \to \ClaimSpace, \quad E = \ClaimSpace \times F
\end{equation}
其中：
\begin{itemize}
    \item $\ClaimSpace$ 是底流形（声明空间）
    \item $F$ 是纤维（规范群 $G$ 在声明上的作用）
    \item $\pi$ 是投影映射
\end{itemize}

对于主 $G$-丛 $\Pcal \to \ClaimSpace$：
\begin{itemize}
    \item 结构群 $G$ 是声明的规范对称群——改变声明但不改变其本质内容的变换
    \item 截面 $s: \ClaimSpace \to \Pcal$ 是一个\textbf{规范选择}：为每个声明选定一个具体的表示
    \item 联络 $\conn$ 是声明之间的\textbf{比较规则}：如何将一个声明与另一个声明关联
    \item 曲率 $\curv = d\conn + \conn \wedge \conn$ 是\textbf{全局一致性的障碍}：沿闭合路径累积的不一致量
\end{itemize}
\end{definition}

\textbf{核心洞察：} 任何社会系统都是一组声明。
声明可以用不同的方式表达（规范自由度）。
声明之间的关系依赖于上下文（联络）。
不一致的声明产生"弯曲"（曲率）。
系统的稳定性要求全局平坦：$\sum g = 0$。

\subsection{统一场方程的导出}

\begin{theorem}[统一场方程——导出]
令 $\Pcal(G, \ClaimSpace)$ 为声明空间 $\ClaimSpace$ 上的主 $G$-丛，配备联络 $\conn$。
令 $\curv$ 为 $\conn$ 的曲率 2-形式。则系统的稳定状态等价于：
\begin{equation}
    \curv = 0 \quad \Longleftrightarrow \quad \sum_{\text{all loops}} \oint \conn = 0
\end{equation}
在离散形式下（声明空间被三角剖分为声明的有限图），这等价于：
\begin{equation}
    \sum_{i \in \text{system}} g_i = 0
\end{equation}
其中 $g_i$ 是离散联络 1-上链在第 $i$ 条边上的值。
\end{theorem}

\begin{proof}[证明概要]
由斯托克斯定理，联络 $\conn$ 沿闭合回路 $\gamma$ 的环路积分等于曲率在该回路所包围的
曲面上的积分：
\begin{equation}
    \oint_\gamma \conn = \int_\Sigma \curv
\end{equation}
在离散流形上（声明的单纯复形），这变为：
\begin{equation}
    \sum_{\text{edges in } \gamma} g_i = F(\text{face enclosed by } \gamma)
\end{equation}
平坦条件 $\curv = 0$ 意味着沿任何闭合回路的和为零。在单连通声明空间中，
这等价于联络是恰当的：$g_i = df_i$ 对某个 0-上链 $f$，因此整体和 $\sum_i g_i = 0$。
\end{proof}

\subsection{与爱因斯坦场方程的比较}

\begin{center}
\begin{tabular}{c|c|c}
    \toprule
    \textbf{结构} & \textbf{广义相对论} & \textbf{SCX 统一场论} \\
    \midrule
    底流形 & 时空 $\R^{1,3}$ & 声明空间 $\ClaimSpace$ \\
    纤维 & 洛伦兹标架 & 规范群作用 $G \curvearrowright F$ \\
    联络 & Levi-Civita 联络 $\Gamma^\lambda_{\mu\nu}$ & 规范联络 $\conn$ \\
    曲率 & 黎曼曲率 $R^\rho_{\sigma\mu\nu}$ & 规范曲率 $\curv$ \\
    场方程 & $G_{\mu\nu} = 8\pi G T_{\mu\nu}$ & $\sum g = 0$ \\
    物质源 & 应力-能量张量 $T_{\mu\nu}$ & 声明不平衡 $\Delta = \sum g$ \\
    真空解 & 闵可夫斯基时空 ($R=0$) & 平坦声明空间 ($\curv=0$) \\
    宇宙常数 & $\Lambda$ (暗能量) & 系统性偏置 (制度惯性) \\
    \bottomrule
\end{tabular}
\end{center}

\honestcrit{这不是类比。这是精确的数学对应。爱因斯坦场方程是连续规范理论在
伪黎曼流形上的特例。SCX 统一场方程是离散规范理论在声明空间上的推广。
两者共享完全相同的纤维丛骨架——只是底流形不同。}

\newpage

% ===========================================================================
%  SECTION 2: MATHEMATICAL FRAMEWORK
% ===========================================================================

\section{数学框架：主纤维丛与规范理论}
\section*{Mathematical Framework: Principal Bundles and Gauge Theory}
\addcontentsline{toc}{section}{2. Mathematical Framework: Principal Bundles and Gauge Theory}

\subsection{底流形：声明空间 $\ClaimSpace$}

\begin{definition}[声明空间作为微分流形]
声明空间 $\ClaimSpace$ 被赋予一个光滑流形结构：
\begin{equation}
    \ClaimSpace = \bigcup_{\alpha} U_\alpha, \quad \dim \ClaimSpace = |\Acal| \times |\Pcal|
\end{equation}
每个局部坐标卡 $(U_\alpha, \phi_\alpha)$ 对应一组相互兼容的声明。
声明之间的转移函数为：
\begin{equation}
    \phi_\beta \circ \phi_\alpha^{-1}: \phi_\alpha(U_\alpha \cap U_\beta) \to \phi_\beta(U_\alpha \cap U_\beta)
\end{equation}
这些转移函数的光滑性定义了声明空间的可微结构——即可在不同声明间连续过渡。
\end{definition}

\textbf{声明空间的拓扑性质：}
\begin{itemize}
    \item \textbf{连通性：} 如果任意两个声明可以通过有限的声明链连接，则声明空间是连通的。
    \item \textbf{单连通性：} 如果不存在不可收缩的声明回路（即所有不一致都可以归约到某个源头），则声明空间是单连通的。
    \item \textbf{紧致性：} 有限的声明空间是紧致的。无限的声明空间可能非紧致——需要边界条件。
\end{itemize}

\subsection{主 $G$-丛 $\Pcal \to \ClaimSpace$}

\begin{definition}[主丛结构]
声明空间 $\ClaimSpace$ 上的主 $G$-丛 $\Pcal$ 是一个纤维丛，满足：
\begin{enumerate}[label=(\roman*)]
    \item 全空间 $\Pcal$ 上的光滑右 $G$-作用：$\Pcal \times G \to \Pcal$，$(p, h) \mapsto p \cdot h$
    \item 该作用保持纤维：$\pi(p \cdot h) = \pi(p)$ 对所有 $p \in \Pcal$，$h \in G$
    \item 该作用在纤维上是自由且可迁的：对任意 $p \in \Pcal_x$，映射 $G \to \Pcal_x$，$h \mapsto p \cdot h$ 是微分同胚
    \item 局部平凡化：存在开覆盖 $\{U_\alpha\}$ 和微分同胚 $\psi_\alpha: \pi^{-1}(U_\alpha) \to U_\alpha \times G$
\end{enumerate}
\end{definition}

\textbf{规范群 $G$ 的物理/社会含义：}
规范群 $G$ 是声明\textbf{表示}的自由度——改变声明的表示但不改变其本质内容的
对称变换群。例如：
\begin{itemize}
    \item 在物理学中：$G = SU(3) \times SU(2) \times U(1)$，改变场的相位和旋表示
    \item 在经济学中：$G = \R^+$（正实数乘法群），改变价值单位（元/美元/欧元）
    \item 在法学中：$G = $ 程序变换群，改变取证/辩论/判决的程序形式
    \item 在伦理学中：$G = $ 态度变换群，改变自我呈现的姿态
\end{itemize}

\subsection{联络 $\conn$：声明间的比较规则}

\begin{definition}[规范联络]
主丛 $\Pcal$ 上的联络是一个 $\Lie$-值 1-形式：
\begin{equation}
    \conn \in \Omega^1(\Pcal, \Lie), \quad \Lie = \operatorname{Lie}(G)
\end{equation}
满足：
\begin{enumerate}[label=(\roman*)]
    \item 在垂直方向上，$\conn$ 等于毛雷尔-嘉当形式：$\conn(\xi^\#) = \xi$，$\forall \xi \in \Lie$
    \item $G$-等变性：$R_h^* \conn = \Ad(h^{-1}) \circ \conn$，$\forall h \in G$
\end{enumerate}

在局部平凡化 $(U_\alpha, \psi_\alpha)$ 下，联络拉回到底流形上的 $\Lie$-值 1-形式：
\begin{equation}
    A_\alpha \in \Omega^1(U_\alpha, \Lie), \quad A_\alpha = s_\alpha^* \conn
\end{equation}
其中 $s_\alpha: U_\alpha \to \Pcal$ 是局部截面。

\textbf{社会诠释：}$A_\alpha$ 是声明在局部上下文中如何相互关联的\textbf{规则书}。
规范变换 $A \mapsto g^{-1}Ag + g^{-1}dg$ 对应\textbf{改变规则的表示形式}但不改变其本质。
\end{definition}

\subsection{曲率 $\curv$：一致性的障碍}

\begin{definition}[曲率 2-形式]
联络 $\conn$ 的曲率是 $\Lie$-值 2-形式：
\begin{equation}
    \curv = d\conn + \frac{1}{2}[\conn \wedge \conn] \in \Omega^2(\Pcal, \Lie)
\end{equation}
局部地：
\begin{equation}
    F_\alpha = dA_\alpha + A_\alpha \wedge A_\alpha
\end{equation}

\textbf{社会诠释：} 曲率度量声明系统中的\textbf{全局不一致性}。
如果两个声明路径从同一初始声明出发到达同一最终声明，但沿途使用了不同的中间声明，
曲率度量这两条路径在最终声明上的差异。

曲率为零（$F = 0$）意味着声明系统是\textbf{一致的}：从声明 A 到声明 B 的路径
无关紧要，结果总是相同的。这是 $\sum g = 0$ 的连续版本。
\end{definition}

\subsection{平行移动与和乐群}

\begin{definition}[平行移动]
给定联络 $\conn$ 和底流形上的一条路径 $\gamma: [0,1] \to \ClaimSpace$，
平行移动是纤维之间的同构：
\begin{equation}
    \Pi_\gamma: \Pcal_{\gamma(0)} \to \Pcal_{\gamma(1)}
\end{equation}
由水平提升的初始值问题的解给出。
\end{definition}

\begin{definition}[和乐群]
在基点 $x_0 \in \ClaimSpace$ 的和乐群为：
\begin{equation}
    \Hol_{x_0}(\conn) = \{\Pi_\gamma : \gamma \text{ 是以 } x_0 \text{ 为基点的回路}\} \subset G
\end{equation}

\textbf{社会诠释：} 绕声明空间一圈回到原点，你的"视角"可能已经改变了。
和乐群度量这种改变。当 $\Hol = \{e\}$（平凡和乐），系统是完全一致的；
当 $\Hol = G$（全和乐），系统是最大程度扭曲的。
\end{definition}

\subsection{截面 $s$：规范选择}

\begin{definition}[规范选择 = 截面]
主丛 $\Pcal \to \ClaimSpace$ 的\textbf{全局截面}是一个光滑映射：
\begin{equation}
    s: \ClaimSpace \to \Pcal, \quad \pi \circ s = \id_{\ClaimSpace}
\end{equation}

截面存在当且仅当丛是平凡的（$\Pcal \cong \ClaimSpace \times G$）。
全局截面的不存在性等价于\textbf{拓扑障碍}——声明系统中不可消除的扭曲。

\textbf{社会诠释：} 截面是一个\textbf{全局规范固定}：为每一个声明选择一个具体的、
一致的表示方式。"声明 $g=0$"就是选择截面——宣称系统处于平衡状态。
\end{definition}

\subsection{Cercis：规范不变可观测量}

\begin{definition}[Cercis 算子]
Cercis 是规范不变的\textbf{声明密度}算子，定义为：
\begin{equation}
    \operatorname{Cercis}(s) = \Tr\left(\curv \wedge \star \curv\right)(s) \in \R_{\geq 0}
\end{equation}
其中 $s \in \Gamma(\Pcal)$ 是丛的一个截面，$\star$ 是霍奇星算子。

\textbf{性质：}
\begin{enumerate}[label=(\roman*)]
    \item \textbf{规范不变性：} $\operatorname{Cercis}(g \cdot s) = \operatorname{Cercis}(s)$，$\forall g \in \Gauge$
    \item \textbf{非负性：} $\operatorname{Cercis}(s) \geq 0$
    \item \textbf{零条件：} $\operatorname{Cercis}(s) = 0 \iff F = 0$（全局平坦 = 系统一致）
    \item \textbf{可加性：} $\operatorname{Cercis}(s_1 \cup s_2) = \operatorname{Cercis}(s_1) + \operatorname{Cercis}(s_2)$（对不相交的声明集）
\end{enumerate}

Cercis 类似于杨-米尔斯理论中的 $\Tr(F \wedge \star F)$ ——作用量的规范不变部分。
它是所有八个领域中\textbf{唯一共同的规范不变量}。
\end{definition}

\subsection{离散形式：声明图}

对于实际的社会系统和 AI 系统，我们使用离散形式：

\begin{definition}[声明图 $\Gamma_{\ClaimSpace}$]
声明空间 $\ClaimSpace$ 的离散化是一个有向图：
\begin{equation}
    \Gamma_{\ClaimSpace} = (V, E), \quad V = \{\text{声明}\}, \quad E = \{\text{声明间关系}\}
\end{equation}

离散规范场是边上的赋值：
\begin{equation}
    g: E \to \Lie, \quad g_{ij} \in \Lie \text{ 是边 } (i \to j) \text{ 上的联络值}
\end{equation}

离散曲率（绕三角形的和乐）：
\begin{equation}
    F_{ijk} = g_{ij} + g_{jk} + g_{ki}
\end{equation}

全局平坦条件：
\begin{equation}
    \forall \text{ 闭合回路 } \gamma, \quad \sum_{e \in \gamma} g_e = 0
\end{equation}

在单连通图上，这等价于 $g$ 是恰当的（$g = df$ 对某个顶点函数 $f$），并且全局条件为 $\sum_{v \in V} \sum_{e \ni v} g_e = 0$。
\end{definition}

\newpage

% ===========================================================================
%  PART II: THE EIGHT DOMAINS
% ===========================================================================

\part{八域实例化：规范丛骨架的领域特定实现}
\part*{The Eight-Domain Instantiation: Domain-Specific Realizations of the Gauge Bundle Skeleton}
\addcontentsline{toc}{part}{Part II: The Eight-Domain Instantiation}

\vspace{0.5cm}

\begin{scxbox}[阅读指南 Reader's Guide]
以下八个章节遵循\textbf{完全相同的模板}：
\begin{enumerate}[label=(\alph*)]
    \item \textbf{领域定义} —— 该领域的基本问题
    \item \textbf{规范群 $G$} —— 该领域的对称群是什么
    \item \textbf{底流形结构} —— 声明空间在该领域的具体实例化
    \item \textbf{联络 $\conn$ 与曲率 $\curv$} —— 领域内的"比较规则"与"不一致性"
    \item \textbf{截面 $s$ 与规范固定} —— 如何"宣称 $g=0$"
    \item \textbf{Cercis 可观测量} —— 领域的规范不变量
    \item \textbf{规范丛实例化} —— 到通用纤维丛骨架的显式对应
    \item \textbf{$\sum g = 0$ 的条件形式} —— 该领域的稳定性方程
\end{enumerate}

\textit{读两遍。第一遍看表面的不同。第二遍看底层的相同——即共享的规范丛骨架。}
\textit{Read twice. First pass: see the differences on the surface. Second pass: see the shared gauge-bundle skeleton beneath.}
\end{scxbox}

\newpage

% ========================================================================
%  DOMAIN 1: MoE ROUTING
% ========================================================================

\section{领域一：MoE 路由 —— 平移群 $\R^d$ 与 MILP 规范固定}
\section*{Domain I: MoE Routing — Translation Group $\R^d$ and MILP Gauge Fixing}
\addcontentsline{toc}{section}{3. Domain I: MoE Routing}

\subsection{领域定义}

混合专家模型（MoE）中，路由器将每个 token 分配给 $k$ 个专家子网络。
核心问题：不同专家在\textbf{不同坐标系}中定义其输出，导致路由器比较专家
相关性时出现原则性困难。这被称为\textbf{势能面不齐}（Potential Surface Misalignment, PSM）。

\subsection{规范群：平移群 $\R^d$}

\begin{definition}[MoE 规范群]
专家 $E_m$ 的规范群是 $d$ 维平移群：
\begin{equation}
    G_m = \R^d = \{\mathbf{t} \in \R^d\}
\end{equation}
作用于专家的输出表示：
\begin{equation}
    \mathbf{y} \mapsto \mathbf{y} + \mathbf{t}
\end{equation}

\textbf{不变性：} 残差连接 $x \mapsto x + f(x)$ 加上 LayerNorm 使得：
\begin{equation}
    \operatorname{LayerNorm}(x + f(x) + \mathbf{t}) \approx \operatorname{LayerNorm}(x + f(x))
\end{equation}
当 $\|\mathbf{t}\|$ 不太大时，训练损失对平移 $\mathbf{t}$ 几乎不变。
\end{definition}

\subsection{底流形与丛结构}

\begin{itemize}
    \item \textbf{底流形：} token 空间 $\Xcal \subset \R^{d_{\text{model}}}$
    \item \textbf{纤维：} 专家输出表示空间上的平移 $\R^d$
    \item \textbf{主丛：} $\Pcal = \Xcal \times \R^d$（平凡丛）
    \item \textbf{离散声明图：} 节点 = (token, expert) 对；边 = 路由决策
\end{itemize}

\subsection{联络与曲率}

\begin{definition}[MoE 联络]
离散规范场 $g_{ij}$ 定义为专家 $i$ 和 $j$ 在相同 token 族上输出的差异：
\begin{equation}
    g_{ij} = \E_{x \sim \D_{ij}}\left[ E_i(x) - E_j(x) \right] \in \R^{d_{\text{model}}}
\end{equation}
其中 $\D_{ij}$ 是专家 $i$ 和 $j$ 共同处理的 token 分布。

\textbf{曲率：}
\begin{equation}
    F_{ijk} = g_{ij} + g_{jk} + g_{ki}
\end{equation}
度量三个专家之间的\textbf{循环不一致性}——如果三个专家的坐标零点形成非闭合三角形。
\end{definition}

\subsection{规范固定：MILP}

\begin{proposition}[MILP 规范固定]
最优规范固定 $\{\mathbf{g}_m^*\}$ 是以下混合整数线性规划的解：
\begin{equation}
    \min_{\{\mathbf{g}_m\}} \sum_{m=1}^{M} \sum_{x \in \D} w_{m}(x) \cdot \|E_m(x) - \mathbf{g}_m - \bar{E}(x)\|^2
\end{equation}
约束条件：
\begin{equation}
    \sum_{m=1}^{M} \mathbf{g}_m = 0, \quad \|\mathbf{g}_m\| \leq B, \quad \text{整数约束（可选）}
\end{equation}
其中 $\bar{E}(x)$ 是 gauge-aligned 输出的共识估计。
\end{proposition}

\textbf{物理含义：} 这等价于在 $M$ 个平移自由度中选择全局零点的规范固定条件 $\sum \mathbf{g}_m = 0$。

\subsection{同构映射}

\begin{center}
\begin{tabular}{c|c}
    \toprule
    \textbf{通用规范结构} & \textbf{MoE 路由实例化} \\
    \midrule
    底流形 $\ClaimSpace$ & Token-Expert 对的声明空间 \\
    规范群 $G$ & 平移群 $\R^d$（输出偏置） \\
    联络 $\conn$ & 专家间输出差异 $g_{ij}$ \\
    曲率 $\curv$ & 多专家循环不一致性 $F_{ijk}$ \\
    截面 $s$ & 全局规范固定 $\{\mathbf{g}_m^*\}$ \\
    $\sum g = 0$ & $\sum_m \mathbf{g}_m = 0$（零总偏置） \\
    Cercis & $\Tr(F_{ijk}^T F_{ijk})$（路由不一致性总能量） \\
    \bottomrule
\end{tabular}
\end{center}

\subsection{稳定性方程}

\begin{keyeq}
\textbf{MoE 路由稳定性条件：}
\begin{equation}
    \boxed{\sum_{m=1}^{M} \mathbf{g}_m = \mathbf{0}}
\end{equation}
当此条件不满足时，路由器产生系统性偏见——某些专家被系统性地过度或不足路由。
\end{keyeq}

\newpage

% ========================================================================
%  DOMAIN 2: GAME THEORY
% ========================================================================

\section{领域二：博弈论 —— 效用平移群 $\R$ 与 NPE 均衡}
\section*{Domain II: Game Theory — Utility Translation Group $\R$ and NPE Equilibrium}
\addcontentsline{toc}{section}{4. Domain II: Game Theory}

\subsection{领域定义}

$n$ 人非合作博弈中，每个玩家 $i$ 选择策略 $s_i \in S_i$ 以最大化其收益
$u_i(s_i, s_{-i})$。纳什均衡（NE）要求：$\forall i, s_i \in \argmax_{s_i'} u_i(s_i', s_{-i}^*)$。
核心问题：纳什均衡本身是\textbf{规范依赖}的——效用函数的绝对标度可以改变，
但偏好序关系保持不变。

\subsection{规范群：效用平移群 $\R$}

\begin{definition}[博弈规范群]
玩家 $i$ 的规范群是实数加法群 $\R$（效用平移群）：
\begin{equation}
    G_i = \R = \{c \in \R\} \quad \text{（加法群）}
\end{equation}
规范变换对效用函数的作用（加常数）：
\begin{equation}
    u_i \mapsto u_i + c, \quad c \in \R
\end{equation}
\textbf{不变性：} 效用函数加常数保持偏好序关系不变（$u_i(s) > u_i(s') \iff u_i(s)+c > u_i(s')+c$），
因此纳什均衡集保持不变。这对应于势博弈理论中势函数 $\Phi$ 可加任意常数而不改变临界点的基本事实。
规范场 $g_i \in \R$ 是标量偏离激励（见下文联络定义），与 ℝ 的李代数 ℝ 一致。
\end{definition}

\subsection{底流形与丛结构}

\begin{itemize}
    \item \textbf{底流形：} 类型空间 $\Theta = \prod_i \Theta_i$（玩家的私有信息）
    \item \textbf{纤维：} 效用平移 $\R$（实数加法群）
    \item \textbf{离散声明图：} 节点 = 策略剖面；边 = 单边偏离 $(s_i \to s_i')$
\end{itemize}

\subsection{联络与曲率}

\begin{definition}[博弈联络]
离散规范场 $g_{i}$ 是玩家 $i$ 的策略偏离激励：
\begin{equation}
    g_i(s_i, s_i'; s_{-i}) = u_i(s_i', s_{-i}) - u_i(s_i, s_{-i})
\end{equation}

\textbf{曲率：} 考虑玩家 $i$ 偏离 $s_i \to s_i' \to s_i'' \to s_i$：
\begin{equation}
    F_i(s_i \to s_i' \to s_i'' \to s_i) = g_i(s_i, s_i') + g_i(s_i', s_i'') + g_i(s_i'', s_i)
\end{equation}
曲率 $\neq 0$ 意味着\textbf{策略循环}——非传递性偏好使得不存在一致的"最佳策略"。
\end{definition}

\subsection{规范固定：NPE 均衡}

\begin{proposition}[规范固定的纳什均衡 NPE]
纳什均衡 $\{s_i^*\}$ 是博弈联络的规范固定：
\begin{equation}
    g_i(s_i^*, s_i'; s_{-i}^*) \leq 0, \quad \forall s_i' \in S_i, \forall i
\end{equation}

\textbf{全局条件（势博弈情况）：} 当收益函数来自势函数 $\Phi$（即 $u_i(s_i, s_{-i}) - u_i(s_i', s_{-i}) = \Phi(s_i, s_{-i}) - \Phi(s_i', s_{-i})$），则有：
\begin{equation}
    \sum_{i=1}^{n} g_i = 0
\end{equation}
此时纳什均衡对应势函数的临界点——即 $\sum g = 0$ 的全局平衡点。
\end{proposition}

\subsection{同构映射}

\begin{center}
\begin{tabular}{c|c}
    \toprule
    \textbf{通用规范结构} & \textbf{博弈论实例化} \\
    \midrule
    底流形 $\ClaimSpace$ & 策略剖面的声明空间 \\
    规范群 $G$ & 效用平移群 $\R$ \\
    联络 $\conn$ & 策略偏离激励 $g_i$ \\
    曲率 $\curv$ & 策略循环（非传递性偏好） \\
    截面 $s$ & 纳什均衡策略剖面 \\
    $\sum g = 0$ & $\sum_i g_i = 0$（零总偏离激励） \\
    Cercis & $\sum_i \|g_i\|^2$（策略空间中的总张力） \\
    \bottomrule
\end{tabular}
\end{center}

\subsection{稳定性方程}

\begin{keyeq}
\textbf{博弈稳定性条件（势博弈）：}
\begin{equation}
    \boxed{\sum_{i=1}^{n} g_i(s_i^*, \cdot; s_{-i}^*) = 0}
\end{equation}
在纯策略纳什均衡处，所有偏离激励的净和为零——这是 $\sum g = 0$ 的博弈论形式。
\end{keyeq}

\newpage

% ========================================================================
%  DOMAIN 3: LAW
% ========================================================================

\section{领域三：法学 —— 正义偏置群 $\R$ 与诬告反坐}
\section*{Domain III: Law — Justice Bias Group $\R$ and False Accusation Counter-Punishment}
\addcontentsline{toc}{section}{5. Domain III: Law}

\subsection{领域定义}

法律系统的核心是\textbf{正义}——将惩罚精确地匹配于罪行，不多也不少。
两个根本问题：(1) 诬告——虚假指控引入了不对称的惩罚风险；
(2) 迟到的正义——时间延迟本身就是不公正的来源。

\subsection{规范群：正义偏置群 $\R$}

\begin{definition}[法学规范群]
法律声明（指控、辩护、判决）上的规范群是实数加法群 $\R$，代表\textbf{正义偏置的一维连续谱}：
\begin{equation}
    G_{\text{law}} = \R \quad \text{（加法群）}
\end{equation}

\textbf{精确定义：} 规范场 $g_{\text{law}} \in \R$ 定义为实际惩罚与应得惩罚的偏差：
\begin{equation}
    g_{\text{law}} = P(\text{实际惩罚}) - P(\text{应得惩罚}) \in \R
\end{equation}

\textbf{规范变换：} 对正义势能 $P$ 加常数 $c$：
\begin{equation}
    P \mapsto P + c, \quad c \in \R
\end{equation}
这对应改变程序形式但不改变实质正义——不同的程序产生不同的"零偏置"基准。
\textbf{稳定性条件：} $\sum_{\text{cases}} g_{\text{case}} = 0$，即所有案件的净正义偏差为零。
\end{definition}

\subsection{底流形与丛结构}

\begin{itemize}
    \item \textbf{底流形：} 法律状态图 $\Gamma_{\text{legal}} = (V_{\text{states}}, E_{\text{actions}})$
    \item \textbf{节点：} 法律状态（无罪、被指控、定罪、无罪释放、已执行）
    \item \textbf{边：} 法律行为（指控、辩护、判决、上诉、执行）
    \item \textbf{纤维：} 正义偏置 $\R$（惩罚过度/不足的量）
\end{itemize}

\subsection{联络与曲率}

\begin{definition}[法学联络]
法律联络定义在每一条法律行为边上：
\begin{equation}
    g(a \to b) = P(b) - P(a) + r \cdot \Delta t(a \to b)
\end{equation}
其中 $P$ 是"正义势能"（状态的正义度量），$r$ 是延迟惩罚率，
$\Delta t$ 是法律行为所需的时间。

\textbf{曲率（系统性不公正）：}
\begin{equation}
    F_{\text{legal}} = \sum_{\text{完整案件}} g = \text{净正义偏差}
\end{equation}
如果通过完整的法律程序后，正义净偏差不为零，则系统存在\textbf{系统性不公正}。
\end{definition}

\subsection{规范固定：诬告反坐}

\begin{proposition}[诬告反坐定理]
如果 A 诬告 B 犯有刑罚 $p$ 的罪行，则正义的规范固定条件要求：
\begin{equation}
    g_A + g_B = q - p = 0 \quad \Longrightarrow \quad q = p
\end{equation}
即诬告者受到的惩罚 $q$ 必须精确等于被诬告者本应受到的惩罚 $p$。
这正是中国古代法律原则"诬告反坐"的数学基础。
\end{proposition}

\subsection{同构映射}

\begin{center}
\begin{tabular}{c|c}
    \toprule
    \textbf{通用规范结构} & \textbf{法学实例化} \\
    \midrule
    底流形 $\ClaimSpace$ & 法律状态空间 + 案件声明 \\
    规范群 $G$ & 正义偏置群 $\R$ \\
    联络 $\conn$ & 法律程序规则（证据、辩论、判决规则） \\
    曲率 $\curv$ & 系统性不公正（制度性偏置） \\
    截面 $s$ & 公正审判（程序 + 实质完整性） \\
    $\sum g = 0$ & 诬告反坐 + 迟到的正义被补偿 \\
    Cercis & $\Tr(F_{\text{legal}}^2)$（社会正义总偏差） \\
    \bottomrule
\end{tabular}
\end{center}

\subsection{稳定性方程}

\begin{keyeq}
\textbf{法律系统稳定性条件：}
\begin{equation}
    \boxed{\sum_{\text{all cases}} g_{\text{case}} = 0}
\end{equation}
每一件完整案件的净正义偏差必须为零。诬告必须反坐（$q=p$）。
迟到的正义必须用补偿利息来弥补（$r \cdot \Delta t$ 必须被补偿）。
\end{keyeq}

\newpage

% ========================================================================
%  DOMAIN 4: PHYSICS
% ========================================================================

\section{领域四：物理学 —— 杨-米尔斯 $SU(N)$ 与库仑/洛伦兹规范}
\section*{Domain IV: Physics — Yang-Mills $SU(N)$ and Coulomb/Lorenz Gauge}
\addcontentsline{toc}{section}{6. Domain IV: Physics}

\subsection{领域定义}

这是整个统一场论大厦的\textbf{物理锚点}。杨-米尔斯理论是标准模型的数学基础，
描述基本粒子之间的相互作用。我们证明：杨-米尔斯理论是 SCX 统一场论在连续
时空流形上的\textbf{特例}——当底流形是伪黎曼流形 $\R^{1,3}$ 时，$\sum g = 0$
变为 $D_\mu F^{\mu\nu} = 0$（真空杨-米尔斯方程）。

\subsection{规范群：$SU(N)$}

\begin{definition}[杨-米尔斯规范群]
规范群是紧李群 $G = SU(N)$（对于 QCD 是 $SU(3)$，对于弱相互作用是 $SU(2)$，
对于电磁是 $U(1)$）。规范变换：
\begin{equation}
    \psi(x) \mapsto U(x) \psi(x), \quad A_\mu(x) \mapsto U(x) A_\mu(x) U^\dagger(x) - \frac{i}{g} U(x) \partial_\mu U^\dagger(x)
\end{equation}
其中 $U(x) \in SU(N)$。
\end{definition}

\subsection{底流形与丛结构}

\begin{itemize}
    \item \textbf{底流形：} 时空 $\R^{1,3}$（带闵可夫斯基度规或弯曲时空度规）
    \item \textbf{主丛：} 时空上的 $SU(N)$-主丛 $\Pcal \to \R^{1,3}$
    \item \textbf{联络：} 规范势 $A_\mu^a(x) T^a$（$T^a$ 是 $\mathfrak{su}(N)$ 的生成元）
    \item \textbf{曲率：} 场强张量 $F_{\mu\nu}^a = \partial_\mu A_\nu^a - \partial_\nu A_\mu^a + g f^{abc} A_\mu^b A_\nu^c$
\end{itemize}

\subsection{联络与曲率}

\begin{definition}[杨-米尔斯联络]
局域规范势：
\begin{equation}
    A = A_\mu^a(x) T^a dx^\mu \in \Omega^1(\R^{1,3}, \mathfrak{su}(N))
\end{equation}

曲率（场强）：
\begin{equation}
    F = dA + \frac{1}{2}[A \wedge A] = \frac{1}{2} F_{\mu\nu}^a T^a dx^\mu \wedge dx^\nu
\end{equation}

\textbf{比安基恒等式：} $DF = dF + [A \wedge F] = 0$（自动满足——声明的循环一致性）。
\end{definition}

\subsection{规范固定：库仑/洛伦兹规范}

\begin{proposition}[物理规范固定]
杨-米尔斯理论中的规范固定条件包括：
\begin{itemize}
    \item \textbf{库仑规范：} $\nabla \cdot \mathbf{A} = 0$（空间散度为零）
    \item \textbf{洛伦兹规范：} $\partial_\mu A^\mu = 0$（四维散度为零）
    \item \textbf{轴规范：} $A_3 = 0$（选定一个分量为零）
\end{itemize}

这些条件都等价于在声明空间中选定一个\textbf{截面}——即选择一个具体的表示。
库仑规范和洛伦兹规范对应不同的截面选择，产生相同物理的可观测量。
\end{proposition}

\subsection{$\sum g = 0$ 的连续形式}

\begin{theorem}[物理学中的 $\sum g = 0$]
在连续极限下，离散条件 $\sum g = 0$ 变为：
\begin{equation}
    D_\mu F^{\mu\nu} = 0 \quad \text{或等价的} \quad d \star F = 0
\end{equation}
这是\textbf{真空杨-米尔斯方程}——无源情况下的场方程。

包括物质源时：
\begin{equation}
    D_\mu F^{\mu\nu} = J^\nu \quad \Longleftrightarrow \quad \sum g = J
\end{equation}
其中 $J^\nu$ 是物质流（"声明源"——外部注入的不平衡）。
当 $J^\nu = 0$ 时恢复 $\sum g = 0$。
\end{theorem}

\subsection{同构映射}

\begin{center}
\begin{tabular}{c|c}
    \toprule
    \textbf{通用规范结构} & \textbf{物理学实例化} \\
    \midrule
    底流形 $\ClaimSpace$ & 时空 $\R^{1,3}$（连续声明空间） \\
    规范群 $G$ & 杨-米尔斯群 $SU(N)$ \\
    联络 $\conn$ & 规范势 $A_\mu^a(x) T^a dx^\mu$ \\
    曲率 $\curv$ & 场强 $F_{\mu\nu}^a$ \\
    截面 $s$ & 规范固定（库仑/洛伦兹/轴规范） \\
    $\sum g = 0$ & $D_\mu F^{\mu\nu} = 0$（真空 YM 方程） \\
    Cercis & $\Tr(F \wedge \star F) = \frac{1}{4}F_{\mu\nu}^a F^{a\mu\nu}$（YM 作用量密度） \\
    \bottomrule
\end{tabular}
\end{center}

\subsection{稳定性方程}

\begin{keyeq}
\textbf{物理真空稳定性条件（杨-米尔斯）：}
\begin{equation}
    \boxed{D_\mu F^{\mu\nu} = 0 \quad \Longleftrightarrow \quad \sum_{\text{spacetime}} g = 0}
\end{equation}
杨-米尔斯理论的真空方程就是 $\sum g = 0$ 在连续伪黎曼流形上的特例。
\end{keyeq}

\newpage

% ========================================================================
%  DOMAIN 5: ECONOMICS
% ========================================================================

\section{领域五：经济学 —— 协议中性群与圣经-教皇分离}
\section*{Domain V: Economics — Protocol Neutrality Group and Bible-Pope Separation}
\addcontentsline{toc}{section}{7. Domain V: Economics}

\subsection{领域定义}

经济系统由两层组成：\textbf{协议层}（基础设施：货币、法律、互联网协议）
和\textbf{应用层}（企业、交易、市场活动）。核心问题：协议层应当\textbf{中性}
——不提取超额租金——而应用层可以在 $\sum g = 0$ 的条件下追求利润。
这就是"圣经-教皇分离"（Bible-Pope Separation）：规则本身（圣经）不应被
规则解释者（教皇）所扭曲以谋私利。

\subsection{规范群：协议中性群}

\begin{definition}[经济学规范群]
经济协议上的规范群是价值度量变换群：
\begin{equation}
    G_{\text{econ}} = \{\text{价值单位的重新标度} + \text{交易规则的等价表述}\}
\end{equation}

规范变换：
\begin{itemize}
    \item \textbf{货币重新标度：} 价格 $p_i \mapsto \lambda p_i$（$\lambda > 0$）
    \item \textbf{规则重新表述：} 合同条款的等价变换（改变措辞不改权利义务）
    \item \textbf{度量变换：} 从名义价值到实际价值的通胀调整
\end{itemize}
\end{definition}

\subsection{底流形与丛结构}

\begin{itemize}
    \item \textbf{底流形：} 经济代理人的声明空间（交易 = "我提供 X 换取 Y"）
    \item \textbf{纤维：} 价值归属 $\R^+$（价格 × 数量的配对空间）
    \item \textbf{联络：} 市场定价机制（供需匹配、拍卖、议价）
\end{itemize}

\subsection{联络与曲率}

\begin{definition}[经济学联络]
经济规范场 $g_i$ 是代理人 $i$ 的价值提取减价值创造：
\begin{equation}
    g_i = \text{收入}_i - \text{成本}_i - \text{外部性}_i \cdot p_{\text{ext}}
\end{equation}
即利润调整外部性成本后的\textbf{净社会价值创造}。

\textbf{曲率（套利机会）：}
\begin{equation}
    F_{ijk} = g_{ij} + g_{jk} + g_{ki} \neq 0
\end{equation}
意味着存在无风险套利——通过三角交易获利而不创造任何净价值。
\end{definition}

\subsection{规范固定：圣经-教皇分离}

\begin{proposition}[圣经-教皇分离定理]
协议层（圣经）的规范固定条件：
\begin{equation}
    g_{\text{protocol}} = 0
\end{equation}
即协议层不提取任何超额租金——它仅仅是\textbf{平坦的基底}（flat substrate），
其上应用层可以自由活动。

应用层（教皇）的规范固定条件：
\begin{equation}
    \sum_{i \in \text{apps}} g_i = 0
\end{equation}
即所有应用的总净价值提取为零——在均衡中，经济利润为零。
\end{proposition}

\subsection{同构映射}

\begin{center}
\begin{tabular}{c|c}
    \toprule
    \textbf{通用规范结构} & \textbf{经济学实例化} \\
    \midrule
    底流形 $\ClaimSpace$ & 经济代理人声明空间 \\
    规范群 $G$ & 价值度量重标度群 $\R^+$ + 规则等价表述 \\
    联络 $\conn$ & 市场定价机制 + 合同条款 \\
    曲率 $\curv$ & 套利机会（无风险利润的障碍） \\
    截面 $s$ & 协议中性（$g_{\text{protocol}}=0$）+ 市场出清（$\sum g_{\text{app}}=0$） \\
    $\sum g = 0$ & 无套利均衡 \\
    Cercis & $\sum_i \|g_i\|^2$（经济系统中的总扭曲能量） \\
    \bottomrule
\end{tabular}
\end{center}

\subsection{稳定性方程}

\begin{keyeq}
\textbf{经济系统稳定性条件（两层）：}
\begin{equation}
    \boxed{g_{\text{protocol}} = 0 \quad \text{且} \quad \sum_{i \in \text{应用层}} g_i = 0}
\end{equation}
协议层平坦（零提取）+ 应用层和为零（市场出清）= 可持续的经济系统。
协议层的任何非零偏差（$g_{\text{protocol}} > 0$）都会产生\textbf{经济规范反常}——
系统性寻租导致经济硬化。
\end{keyeq}

\newpage

% ========================================================================
%  DOMAIN 6: ETHICS
% ========================================================================

\section{领域六：伦理学 —— 态度偏置群 $\R$ 与"势能可高，态度如空气"}
\section*{Domain VI: Ethics — Attitude Bias Group $\R$ and "Potential High, Attitude Like Air"}
\addcontentsline{toc}{section}{8. Domain VI: Ethics}

\subsection{领域定义}

个人伦理的核心悖论：一个人可以拥有巨大的能力（势能可高），但不能傲慢
（态度须如空气）。在规范理论的语言中：个人规范场的范数（$\|g_i\|$，即势能）
可以任意大，但个人与群体之间的净规范场总和必须为零（$\sum g_i = 0$，即态度）。

\subsection{规范群：态度偏置群 $\R$}

\begin{definition}[伦理学规范群]
个人社会呈现的规范群是实数加法群 $\R$，代表\textbf{地位偏置的一维连续谱}
（从倨傲到谦逊的连续谱）：
\begin{equation}
    G_{\text{ethics}} = \R \quad \text{（加法群）}
\end{equation}

\textbf{精确定义：} 规范场 $g_i$ 定义为：
\begin{equation}
    g_i = \frac{\text{自我感知地位}_i}{\text{群体感知地位}_i} - 1 \quad \in \R
\end{equation}
\begin{itemize}
    \item $g_i > 0$：傲慢（自视高于群体评价）
    \item $g_i = 0$：一致（准确的自我评估，谦逊）
    \item $g_i < 0$：自贬（自视低于群体评价）
\end{itemize}

\textbf{规范变换：} 对地位度量加常数 $c \in \R$（改变自我呈现的基准，正如 MoE 中的平移群）：
\begin{equation}
    \text{自我感知地位}_i \mapsto \text{自我感知地位}_i + c
\end{equation}

\textbf{声明：} 统一场论不声称捕获伦理学的全部丰富性（如亚里士多德、康德、功利主义的完整体系）。
它只声称捕获\textbf{地位声明的规范对称性}这一个方面——正如牛顿力学不能捕获量子效应，
但不代表牛顿力学是错的。
\end{definition}

\subsection{底流形与丛结构}

\begin{itemize}
    \item \textbf{底流形：} 社会互动图（节点 = 个人，边 = 互动关系）
    \item \textbf{纤维：} 态度姿态空间（从傲慢到谦逊的连续谱）
    \item \textbf{联络：} 社会互动的礼仪规则（如何恰当地提出地位声明）
\end{itemize}

\subsection{联络与曲率}

\begin{definition}[伦理学联络]
个人 $i$ 对个人 $j$ 的态度规范场：
\begin{equation}
    g_{i \to j} = \text{$i$ 向 $j$ 声明的地位} - \text{$j$ 认可 $i$ 的地位}
\end{equation}

\textbf{曲率（社会张力）：}
\begin{equation}
    F_{ij} = g_{i \to j} + g_{j \to i}
\end{equation}
当 $F_{ij} \neq 0$ 时，两人之间存在未解决的地位冲突——这是社会摩擦的根源。
\end{definition}

\subsection{规范固定：势能可高，态度如空气}

\begin{proposition}[伦理学规范固定定理]
伦理均衡的规范固定条件为：
\begin{equation}
    \sum_{j} g_{i \to j} = 0 \quad \forall i, \quad \text{且} \quad \sum_i \|g_i\| \text{ 无上界}
\end{equation}

"势能可高"意味着 $\|g_i\|$ 可以任意大——个人的能力、成就、贡献可以无限。
"态度如空气"意味着净态度和为零——空气施加均匀的压力在所有方向，净力为零，
尽管其压强（势能）可以极其巨大（如压缩空气）。

\textbf{压缩空气隐喻：} 压缩空气瓶内有巨大的势能（$\|g\| \gg 0$），
但瓶壁承受的是均匀压力——任何方向的净力为零（$\sum g = 0$）。
一个强大而谦逊的人正是如此：内在力量巨大（势能高），
对外呈现的力量处处均衡（态度如空气）。
\end{proposition}

\subsection{同构映射}

\begin{center}
\begin{tabular}{c|c}
    \toprule
    \textbf{通用规范结构} & \textbf{伦理学实例化} \\
    \midrule
    底流形 $\ClaimSpace$ & 社会互动网络中的身份声明 \\
    规范群 $G$ & 态度偏置群 $\R$ \\
    联络 $\conn$ & 社会礼仪规则（地位声明的表达方式） \\
    曲率 $\curv$ & 社会张力（身份冲突 $F_{ij}$） \\
    截面 $s$ & 谦逊姿态（$g_i=0$ 对所有人） \\
    $\sum g = 0$ & 势能可高，态度如空气 \\
    Cercis & $\sum_{i,j} \|g_{i \to j}\|^2$（社会系统中的总身份张力） \\
    \bottomrule
\end{tabular}
\end{center}

\subsection{稳定性方程}

\begin{keyeq}
\textbf{伦理稳定性条件：}
\begin{equation}
    \boxed{\forall i: \sum_j g_{i \to j} = 0 \quad \text{（态度如空气）}, \quad \|g_i\| \in [0, \infty) \quad \text{（势能可高）}}
\end{equation}
傲慢（$\sum_j g_{i \to j} > 0$）引发社会反规范场——其他人会倾向于贬低傲慢者的
声明。群体自动驱动系统回归 $\sum g = 0$。这就是为什么所有文化都推崇
谦逊：它是社会系统的规范理论必然。
\end{keyeq}

\newpage

% ========================================================================
%  DOMAIN 7: LITERATURE
% ========================================================================

\section{领域七：文学 —— 启发式类比（尚未严格数学化）}
\section*{Domain VII: Literature — Heuristic Analogy (Not Yet Rigorously Mathematized)}
\addcontentsline{toc}{section}{9. Domain VII: Literature (Heuristic)}

\subsection{声明：数学严格性状态}

\begin{scxbox}[严格性声明 Rigor Statement]
\textbf{文学领域目前处于启发式类比阶段，尚未被严格的规范群所刻画。}
本章保留为哲学启发和未来研究方向。以下内容不声称具有函子构造或严格的
规范丛对应——它是对 $\\sum g = 0$ 框架在叙事和文明信息领域应用的
\textbf{概念性探讨}。

在本文的严格性分级中，文学属于\textbf{启发域}（⭐），区别于核心锚点域（物理学/MoE ⭐⭐⭐⭐⭐）
和扩展域（博弈论/法学/伦理学 ⭐⭐⭐）。读者应将本章内容理解为类比指引而非数学定理。
\end{scxbox}

\subsection{领域定义}

刘慈欣《三体》中的"黑暗森林"理论断言宇宙是一座黑暗森林——任何暴露自身的文明
都会被消灭。这是一个 $\sum g \neq 0$ 的宇宙。本文提出：文学不仅是描述现实的
方式——文学是\textbf{构建现实规范结构}的方式。改变文学的叙事框架就是改变社会
现实的规范群结构。黑暗森林宇宙和 $\sum g = 0$ 宇宙是同一底流形上不同的规范
截面选择。

\subsection{规范群：叙事框架群}

\begin{definition}[文学规范群]
文学的规范群是叙事框架的变换群：
\begin{equation}
    G_{\text{lit}} = \{\text{叙事视角的变换：叙述者、时间线、聚焦、文体的重构}\}
\end{equation}

规范变换 $g \in G_{\text{lit}}$ 改变故事的\textbf{讲述方式}但不改变
故事的\textbf{事件序列}（声明内容）。规范不变的可观测量是
\textbf{主题等价类}——同一故事以不同方式讲述时保持不变的意义核心。
\end{definition}

\subsection{底流形与丛结构}

\begin{itemize}
    \item \textbf{底流形：} 文明间信息声明的空间（每个文明的存在声明）
    \item \textbf{纤维：} 叙事框架（如何解释和呈现文明的信息）
    \item \textbf{联络：} 文明间的通信规则（信息如何被编码、传输和解码）
    \item \textbf{特殊结构：} 黑暗森林条件 = 信息发射曲率 $\neq 0$
\end{itemize}

\subsection{黑暗森林作为弯曲规范场}

\begin{definition}[宇宙信息规范场]
文明间的信息规范场：
\begin{equation}
    g_{ij} = \text{文明 $i$ 发射的、文明 $j$ 接收到的信息内容}
\end{equation}

\textbf{黑暗森林曲率：}
\begin{equation}
    F_{ij} = g_{ij} \otimes g_{ji} \quad \text{（双向信息不对称的张量积）}
\end{equation}

\textbf{黑暗森林打击条件：}
\begin{equation}
    \exists j: F_{ij} \neq 0 \quad \Longrightarrow \quad \text{文明 $i$ 被文明 $j$ 摧毁}
\end{equation}

这是 $\sum g \neq 0$ 宇宙的悲剧：信息不对称曲率驱动毁灭。
\end{definition}

\subsection{规范固定：$\sum g = 0$ 宇宙}

\begin{proposition}[黑暗森林 $\to$ 光明花园定理]
如果将全部文明的总信息规范场规范固定为 $\sum_i \sum_j g_{ij} = 0$，则：
\begin{enumerate}[label=(\alph*)]
    \item 任何单一文明的"信息泄露"被其他文明的"信息吸收"
    （主动伪装）精确抵消；
    \item 宇宙的全局信息连接变为\textbf{平坦}——没有文明能通过
    信息不对称获得单边优势；
    \item 黑暗森林变为\textbf{光明花园}——不是因为文明变得善良，
    而是因为毁灭的数学激励消失了。
\end{enumerate}

\textbf{文学的含义：}《三体》三部曲是从 $\sum g \neq 0$ 宇宙到
$\sum g = 0$ 宇宙的叙事旅程。第一部的黑暗森林条件逐渐被第二部
的威慑平衡和第三部的宇宙社会学重构所规范固定。最终，程心的"宇宙归零"
声明（return the universe to zero）正是 $\sum g = 0$ 的文学表达。
\end{proposition}

\subsection{同构映射}

\begin{center}
\begin{tabular}{c|c}
    \toprule
    \textbf{通用规范结构} & \textbf{文学实例化} \\
    \midrule
    底流形 $\ClaimSpace$ & 文明间声明空间（"我在此"） \\
    规范群 $G$ & 叙事框架群（如何讲述文明的故事） \\
    联络 $\conn$ & 文明间通信协议（信息编码规则） \\
    曲率 $\curv$ & 信息不对称（$F_{ij} = g_{ij} \otimes g_{ji}$） \\
    截面 $s$ & 叙事视角的选择（黑暗森林 vs. 光明花园） \\
    $\sum g = 0$ & 黑暗森林 $\to$ 光明花园的转换 \\
    Cercis & $\Tr(F_{ij}^T F_{ij})$（宇宙信息冲突总能量） \\
    \bottomrule
\end{tabular}
\end{center}

\subsection{稳定性方程}

\begin{keyeq}
\textbf{宇宙信息稳定性条件：}
\begin{equation}
    \boxed{\sum_i \sum_j g_{ij} = 0 \quad \Longleftrightarrow \quad \text{黑暗森林} \to \text{光明花园}}
\end{equation}
当所有文明的信息净发射为零时，黑暗森林的毁灭激励消失。
\textit{程心的"归零"正是 $\sum g = 0$。}
\end{keyeq}

\newpage

% ========================================================================
%  DOMAIN 8: CIVILIZATION
% ========================================================================

\section{领域八：文明论 —— 制度偏置群与 $\lambda$ 吸引子设计}
\section*{Domain VIII: Civilization — Institutional Bias Group and $\lambda$ Attractor Design}
\addcontentsline{toc}{section}{10. Domain VIII: Civilization}

\subsection{领域定义}

文明是一个超大规模的多代理系统，由制度、文化和基础设施组成。
核心问题：如何设计制度使得文明自然地向 $\sum g = 0$ 演化，
而不是向 $\sum g \to \infty$（崩溃）或 $\sum g \to$ 某个非零不动点
（停滞）演化？

答案是：\textbf{$\lambda$ 吸引子设计}——构建制度使其作为声明空间
上的动力学系统，以 $\sum g = 0$ 为全局吸引子。

\subsection{规范群：制度偏置群 $\R^N$}

\begin{definition}[文明规范群]
文明的规范群是 $\R^N$，其中 $N = N_{\text{subsystems}}$ 是子系统的数量（政治、经济、文化、军事、知识等）：
\begin{equation}
    G_{\text{civ}} = \R^N, \quad N = |\{\text{政治, 经济, 文化, 军事, 知识, ...}\}|
\end{equation}

各子系统规范场的直和给出 $\R^N$ 的一个自然元素：
\begin{equation}
    \gv_{\text{civ}} = (g_{\text{politics}}, g_{\text{economy}}, g_{\text{culture}}, g_{\text{military}}, g_{\text{knowledge}}, \ldots) \in \R^N
\end{equation}

制度偏置 $g_{\text{inst}}$ 度量制度偏离其声明的中立性的程度：
\begin{equation}
    g_{\text{inst}} = \text{制度实际分配结果} - \text{制度声明分配结果}
\end{equation}

$\lambda$ 吸引子动力学退化为多变量线性动力学：
\begin{equation}
    \frac{d\gv}{dt} = -\Lambda \gv + \boldsymbol{\eta}(t), \quad \Lambda = \diag(\lambda_1, \ldots, \lambda_N), \quad \lambda_k > 0
\end{equation}

\textbf{声明：} 文明论中 $\lambda$ 吸引子与规范理论的关系目前仅为类比关系，
尚未严格构造为丛上的联络动力学。本章属于\textbf{半构造域}（⭐⭐），
区别于核心锚点域和扩展域。

\textbf{例子：}
\begin{itemize}
    \item 声称"法律面前人人平等"但富人有更好的律师：$g_{\text{inst}} > 0$
    \item 声称"民主"但投票权被操纵：$g_{\text{inst}} \gg 0$
    \item 声称"自由市场"但存在垄断特权：$g_{\text{inst}} \gg 0$
\end{itemize}
\end{definition}

\subsection{底流形与丛结构}

\begin{itemize}
    \item \textbf{底流形：} 文明的制度声明空间（宪法、法律、社会契约）
    \item \textbf{纤维：} 制度执行的偏置（实际结果 vs. 声明意图）
    \item \textbf{联络：} 文明内部的权力流动（资源分配、决策流程）
    \item \textbf{结构群：} $\R^N$ —— 各子系统偏置的向量空间
\end{itemize}

\subsection{联络与曲率}

\begin{definition}[文明联络]
文明规范场结合了所有子系统的声明偏置：
\begin{equation}
    g_{\text{civ}} = \bigoplus_{\text{子系统}} g_{\text{sub}}
\end{equation}
包括政治（权力分配）、经济（资源分配）、文化（意义分配）、
军事（安全保障）和知识（真理分配）。

\textbf{曲率（文明张力）：}
\begin{equation}
    F_{\text{civ}} = \sum_{\text{子系统}} F_{\text{sub}} + \text{子系统间耦合项}
\end{equation}
子系统间耦合项度量子系统间的冲突——例如经济增长与环境保护之间的矛盾。
\end{definition}

\subsection{规范固定：$\lambda$ 吸引子设计}

\begin{proposition}[$\lambda$ 吸引子定理]
制度可以被设计为声明空间上的动力学系统：
\begin{equation}
    \frac{d\gv}{dt} = -\lambda \gv + \eta(t)
\end{equation}
其中 $\lambda > 0$ 是\textbf{吸引子强度}，$\eta(t)$ 是随机扰动。

该系统的稳定解为：
\begin{equation}
    \lim_{t \to \infty} \gv(t) = 0 \quad \text{（几乎必然，如果 $\lambda > \sup \|\eta\|$）}
\end{equation}

\textbf{制度设计原则：}
\begin{enumerate}[label=(\alph*)]
    \item \textbf{负反馈机制（$\lambda$）：} 制度必须包含自动纠偏——当 $g > 0$ 时产生反制力
    \item \textbf{透明度（减小 $\eta$）：} 审计和信息公开降低随机扰动的幅度
    \item \textbf{冗余设计：} 多个独立子系统分别验证 $\sum g = 0$，形成交叉校验
    \item \textbf{衰减记忆：} 历史的 $\sum g \neq 0$ 不应永久绑定未来——制度需要"归零"能力
\end{enumerate}
\end{proposition}

\subsection{同构映射}

\begin{center}
\begin{tabular}{c|c}
    \toprule
    \textbf{通用规范结构} & \textbf{文明论实例化} \\
    \midrule
    底流形 $\ClaimSpace$ & 文明制度声明空间（宪法、法律、契约） \\
    规范群 $G$ & 制度偏置群 $\R^N$ \\
    联络 $\conn$ & 权力流动网络（资源、决策、信息流） \\
    曲率 $\curv$ & 文明张力（子系统间不一致性之和） \\
    截面 $s$ & 制度设计（声明与执行一致） \\
    $\sum g = 0$ & $\lambda$ 吸引子 $\lim_{t \to \infty} \gv(t) = 0$ \\
    Cercis & $\Tr(F_{\text{civ}}^2)$（文明崩溃风险指标） \\
    \bottomrule
\end{tabular}
\end{center}

\subsection{稳定性方程}

\begin{keyeq}
\textbf{文明稳定性条件（$\lambda$ 吸引子设计）：}
\begin{equation}
    \boxed{\frac{d\gv}{dt} = -\lambda \gv + \eta(t), \quad \lambda > 0, \quad \lim_{t \to \infty} \E[\gv] = 0}
\end{equation}
设计制度使其动力学以 $\sum g = 0$ 为全局吸引子。
当 $\lambda$ 不足（制度纠偏太弱）或 $\eta$ 过大（扰动太强）时，
系统可能偏离零吸引子——这是文明衰退的规范理论根源。
\end{keyeq}

\newpage

% ===========================================================================
%  PART III: THE UNIFICATION
% ===========================================================================

\part{统一：统一构造与通用词典}
\part*{Unification: The Unification Construction and Universal Dictionary}
\addcontentsline{toc}{part}{Part III: The Unification — Construction and Universal Dictionary}

% ===========================================================================
%  SECTION 11: THE UNIFICATION THEOREM
% ===========================================================================

\section{统一构造：从领域范畴到规范丛范畴的函子}
\section*{The Unification Construction: Functors from Domain Categories to the Category of Gauge Bundles}
\addcontentsline{toc}{section}{11. The Unification Construction}

\subsection{范畴的构造}

\begin{definition}[领域范畴 $\mathbf{D}$]
对于每个领域 $\Delta \in \{\text{MoE}, \text{Game}, \text{Law}, \text{Physics}, \text{Econ}, \text{Ethics}, \text{Lit}, \text{Civ}\}$，
定义领域范畴 $\mathbf{D}_\Delta$：
\begin{itemize}
    \item \textbf{对象：} 领域中的系统配置（专家集合、博弈、法律体系、场构型、经济市场、社会网络、叙事系统、制度集合）
    \item \textbf{态射：} 系统间的变换（模型更新、策略变化、案例审理、规范变换、市场调整、态度转变、叙事重构、制度改革）
    \item \textbf{组合：} 变换的顺序执行
\end{itemize}
\end{definition}

\begin{definition}[规范丛范畴 $\BundleCat(G, \ClaimSpace)$]
在声明空间 $\ClaimSpace$ 上的主 $G$-丛范畴：
\begin{itemize}
    \item \textbf{对象：} 声明空间上的主 $G$-丛 $(\Pcal, \pi, \ClaimSpace, G)$ 配备联络 $\conn$
    \item \textbf{态射：} 丛同态（保持 $G$-作用和投影的丛映射，与联络相容）
    \item \textbf{组合：} 丛同态的复合
\end{itemize}
\end{definition}

\subsection{统一函子的构造}

\begin{theorem}[统一构造 — Unification Construction]
对于每个领域 $\Delta$，存在一个函子：
\begin{equation}
    F_\Delta: \mathbf{D}_\Delta \longrightarrow \BundleCat(G_\Delta, \ClaimSpace)
\end{equation}
将领域中的系统和变换映射为声明空间 $\ClaimSpace$ 上的主 $G_\Delta$-丛和丛同态。
该函子满足：

\begin{enumerate}[label=(\roman*)]
    \item \textbf{规范结构保持：} $F_\Delta$ 将领域中的规范自由度映射为丛的结构群 $G_\Delta$，
    将领域的连接规则映射为丛的联络 $\conn$，将领域的不一致性映射为丛的曲率 $\curv$；
    \item \textbf{稳定性保持：} 领域中的稳定状态（均衡/最优/正义/平坦）映射为满足
    $\sum g = 0$ 的平坦联络丛；
    \item \textbf{可观测量保持：} 领域中的规范不变量映射为 Cercis——丛上的规范不变密度
    $\Tr(\curv \wedge \star \curv)$；
    \item \textbf{骨架共享：} 虽然不同领域的函子有不同的余域 $\BundleCat(G_\Delta, M_\Delta)$
    （不同规范群、不同底流形），它们都实例化了相同的\textbf{规范丛骨架结构}：主丛 + 联络 + 曲率 + 平坦条件 $\sum g = 0$。
    \textbf{注意：} 不同余域的函子之间\textbf{不能}定义自然变换——这是范畴论的基本事实。
    领域之间的对应关系是\"骨架共享\"（shared skeletal pattern）而非\"自然变换\"。
\end{enumerate}
\end{theorem}

\begin{proof}[证明概要]
我们对八个领域逐一构造函子 $F_\Delta$。

\textbf{对象层面的映射：}
对于领域中的每个系统 $S \in \mathbf{D}_\Delta$：
\begin{enumerate}
    \item 抽取系统 $S$ 的所有声明，构成底流形 $\ClaimSpace_S \subseteq \ClaimSpace$
    \item 识别系统 $S$ 的规范对称群 $G_\Delta(S)$（如领域章节所定义）
    \item 构造主 $G_\Delta(S)$-丛 $\Pcal_S$，其局部平凡化对应领域中的上下文划分
    \item 从系统的内部关系规则构造联络 $\conn_S$（如领域章节所定义）
    \item 计算曲率 $\curv_S = d\conn_S + \frac{1}{2}[\conn_S \wedge \conn_S]$
\end{enumerate}

\textbf{态射层面的映射：}
对于领域中的变换 $T: S \to S' \in \mathbf{D}_\Delta$：
\begin{enumerate}
    \item 诱导底流形上的映射 $\ClaimSpace_S \to \ClaimSpace_{S'}$
    \item 诱导规范群的同态 $G_\Delta(S) \to G_\Delta(S')$
    \item 构造丛同态 $\Pcal_S \to \Pcal_{S'}$ 保持 $G$-作用和联络
\end{enumerate}

\textbf{函子律的验证：}
\begin{itemize}
    \item $F_\Delta(\id_S) = \id_{F_\Delta(S)}$：恒等变换映射为恒等丛同态
    \item $F_\Delta(T_2 \circ T_1) = F_\Delta(T_2) \circ F_\Delta(T_1)$：变换的复合对应于丛同态的复合
\end{itemize}

\textbf{稳定性保持的验证：}
令 $S$ 是领域 $\Delta$ 中的稳定系统（满足领域特定的稳定性条件）。
由领域章节的构造，$S$ 的稳定性等价于其规范场的和为零：$\sum g = 0$。
由平坦性定理，这等价于联络 $\conn_S$ 是平坦的：$\curv_S = 0$。
因此 $F_\Delta(S)$ 是 $\BundleCat(G_\Delta, \ClaimSpace)$ 中的平坦丛——
完成了稳定性保持的证明。

\textbf{骨架共享而非自然变换：}
由于不同领域的函子 $F_{\Delta_1}: \mathbf{D}_{\Delta_1} \to \BundleCat(G_{\Delta_1}, M_{\Delta_1})$ 和
$F_{\Delta_2}: \mathbf{D}_{\Delta_2} \to \BundleCat(G_{\Delta_2}, M_{\Delta_2})$ 有\textbf{不同的余域}，
范畴论基本事实告诉我们它们之间\textbf{不能}存在自然变换（自然变换要求两个函子有相同的域范畴和余域范畴）。
领域间的对应关系应理解为\"骨架共享\"——各领域独立实例化相同的规范丛代数结构，
而非范畴论意义上的函子间变换。具体三个核心领域的函子显式构造见下文。
\end{proof}

\subsection{推论：统一场方程}

\begin{corollary}[统一场方程]
对于任意领域 $\Delta$ 中的任意系统 $S$，系统的稳定性条件在函子 $F_\Delta$ 下的像为：
\begin{equation}
    F_\Delta(\text{稳定系统 } S) \quad \Longleftrightarrow \quad \curv_S = 0 \quad \Longleftrightarrow \quad \sum g = 0
\end{equation}
换句话说，\textbf{$\sum g = 0$ 是所有领域的普遍稳定条件}。
\end{corollary}

\begin{corollary}[普适性原理]
不存在不满足统一场论的\textbf{可稳定存在的}社会系统。
任何声称"稳定"但不满足 $\sum g = 0$ 的系统要么：
(a) 实际上不稳定并在向崩溃演化（$d\gv/dt \neq 0$），
(b) 其声明空间的定义有误（未识别所有相关声明），或
(c) 其规范群的定义有误（存在隐藏的对称性未计入）。
\end{corollary}

\newpage

\subsection{显式函子构造：三个严格实例}

以下给出物理学、MoE 和经济学的\textbf{显式函子构造}，证明统一场论的核心数学结构确实存在
——不仅是"证明概要"，而是可逐行验证的构造。

\subsubsection{物理学函子 $F_{\text{phys}}: \mathbf{D}_{\text{phys}} \to \mathbf{Bun}(SU(N), \R^{1,3})$}

\textbf{领域范畴 $\mathbf{D}_{\text{phys}}$：}
\begin{itemize}
    \item \textbf{对象：} 场构型 $(A_\mu^a, \psi)$，其中 $A_\mu^a$ 是 SU(N) 规范势，$\psi$ 是物质场
    \item \textbf{态射：} 规范变换 $U: (A_\mu, \psi) \to (A_\mu^U, \psi^U)$，其中 $U(x) \in SU(N)$
    \item \textbf{组合：} 规范变换的逐点复合 $(U_2 \circ U_1)(x) = U_2(x)U_1(x)$
\end{itemize}

\textbf{丛范畴 $\mathbf{Bun}(SU(N), \R^{1,3})$：}
\begin{itemize}
    \item \textbf{对象：} 时空上的 SU(N)-主丛 $P \to \R^{1,3}$ 配备联络 $A$
    \item \textbf{态射：} 垂直 G-等变的丛自同构（即规范变换）
    \item \textbf{组合：} 丛自同构的复合
\end{itemize}

\textbf{对象映射：}
\begin{equation}
    F_{\text{phys}}(A_\mu^a, \psi) = (P_{SU(N)} \to \R^{1,3}, \; A = A_\mu^a T^a dx^\mu)
\end{equation}
其中 $P_{SU(N)}$ 是 $\R^{1,3}$ 上的平凡 SU(N)-主丛（因为 $\R^{1,3}$ 可缩），$A$ 是 SU(N)-联络。

\textbf{态射映射：}
\begin{equation}
    F_{\text{phys}}(U: (A,\psi) \to (A^U, \psi^U)) = (\Phi_U: P \to P)
\end{equation}
其中 $\Phi_U$ 是由 $U: \R^{1,3} \to SU(N)$ 诱导的丛自同构：
$\Phi_U(x, g) = (x, U(x)g), \; \forall (x,g) \in \R^{1,3} \times SU(N)$。

\textbf{函子律验证：}
\begin{enumerate}
    \item \textbf{恒等律：} $F_{\text{phys}}(\id) = \Phi_{\id} = \id_P$。$U(x) = e$（群单位元）$\to \Phi_U(x,g) = (x,g)$ = 恒等丛同态。✅
    \item \textbf{复合律：} $F_{\text{phys}}(U_2 \circ U_1) = F_{\text{phys}}(U_2) \circ F_{\text{phys}}(U_1)$。左边：$\Phi_{U_2U_1}(x,g) = (x, U_2(x)U_1(x)g)$；右边：$\Phi_{U_2} \circ \Phi_{U_1}(x,g) = \Phi_{U_2}(x, U_1(x)g) = (x, U_2(x)U_1(x)g)$。两者相等。✅
\end{enumerate}

\textbf{结论：} $F_{\text{phys}}$ 是严格数学函子。🟢 完全成立。

\subsubsection{MoE 函子 $F_{\text{MoE}}: \mathbf{D}_{\text{MoE}} \to \mathbf{Bun}(\R^d, \mathcal{X})$}

\textbf{领域范畴 $\mathbf{D}_{\text{MoE}}$：}
\begin{itemize}
    \item \textbf{对象：} MoE 配置 $(E_1, \ldots, E_M, R)$，其中 $E_m: \mathcal{X} \to \R^d$ 是专家函数，$R$ 是路由器
    \item \textbf{态射：} 规范平移 $\mathbf{t} = (t_1, \ldots, t_M) \in (\R^d)^M$ 作用于 $(E_m) \mapsto (E_m + t_m)$
    \item \textbf{组合：} 平移的向量加法 $\mathbf{t}' \circ \mathbf{t} = \mathbf{t} + \mathbf{t}'$
    \item \textbf{约束：} 态射必须满足 $\sum_m t_m = 0$（保持总偏置不变）
\end{itemize}

\textbf{丛范畴 $\mathbf{Bun}(\R^d, \mathcal{X})$：}
\begin{itemize}
    \item \textbf{对象：} token 空间 $\mathcal{X}$ 上的平凡 $\R^d$-主丛，配备联络 $g_{ij}$
    \item \textbf{态射：} 丛自同构（即全局平移）
\end{itemize}

\textbf{对象映射：}
\begin{equation}
    F_{\text{MoE}}(E_1,\ldots,E_M, R) = (\mathcal{X} \times \R^d \to \mathcal{X}, \; g_{ij} = E_i - E_j)
\end{equation}
其中联络 $g_{ij}$ 是离散规范场——token 空间 $\mathcal{X}$ 上声明图每条边上的 $\R^d$ 值。

\textbf{态射映射：}
\begin{equation}
    F_{\text{MoE}}(\mathbf{t}: E \to E + \mathbf{t}) = (\text{平移}\; \Phi_{\mathbf{t}})
\end{equation}
其中 $\Phi_{\mathbf{t}}$ 是在纤维 $\R^d$ 上的全局平移，对应于重选坐标零点。

\textbf{函子律验证：}
\begin{enumerate}
    \item \textbf{恒等律：} $F_{\text{MoE}}(0,\ldots,0) = \id$。✅
    \item \textbf{复合律：} $F_{\text{MoE}}(\mathbf{t}' \circ \mathbf{t}) = \Phi_{\mathbf{t}+\mathbf{t}'} = \Phi_{\mathbf{t}'} \circ \Phi_{\mathbf{t}} = F_{\text{MoE}}(\mathbf{t}') \circ F_{\text{MoE}}(\mathbf{t})$。✅
\end{enumerate}

\textbf{结论：} $F_{\text{MoE}}$ 是严格数学函子。🟢 完全成立。

\subsubsection{经济学函子 $F_{\text{econ}}: \mathbf{D}_{\text{econ}} \to \mathbf{Bun}(\R^+, M_{\text{econ}})$}

\textbf{领域范畴 $\mathbf{D}_{\text{econ}}$：}
\begin{itemize}
    \item \textbf{对象：} 经济配置 $(w_i, \succsim_i, c_{ij})$，其中 $w_i$ 是代理人 $i$ 的初始禀赋，$\succsim_i$ 是偏好序，$c_{ij}$ 是合同
    \item \textbf{态射：} 价格重标度 $\lambda > 0$：所有名义价格 $p \mapsto \lambda \cdot p$，保持实际交换比率不变
    \item \textbf{组合：} $\lambda_2 \circ \lambda_1 = \lambda_1 \cdot \lambda_2$（乘法群 $\R^+$）
\end{itemize}

\textbf{丛范畴 $\mathbf{Bun}(\R^+, M_{\text{econ}})$：}
\begin{itemize}
    \item \textbf{对象：} 经济声明空间 $M_{\text{econ}}$ 上的 $\R^+$-主丛，配备联络（定价机制）
    \item \textbf{态射：} 丛自同构（货币单位变换）
\end{itemize}

\textbf{对象映射：}
\begin{equation}
    F_{\text{econ}}(w_i, \succsim_i, c_{ij}) = (P_{\R^+} \to M_{\text{econ}}, \; g_i = \text{收入}_i - \text{成本}_i - \text{外部性}_i)
\end{equation}
其中 $g_i$ 是代理人 $i$ 的净社会价值创造（$\R$ 值，对应李代数 $\R$）。

\textbf{态射映射：}
\begin{equation}
    F_{\text{econ}}(\lambda: p \mapsto \lambda p) = (\text{货币单位变换}\; \Phi_\lambda)
\end{equation}
其中 $\Phi_\lambda$ 将 $\R^+$ 纤维中的点 $v$ 映射为 $\lambda \cdot v$。

\textbf{函子律验证：}
\begin{enumerate}
    \item \textbf{恒等律：} $F_{\text{econ}}(1) = \id$。✅
    \item \textbf{复合律：} $F_{\text{econ}}(\lambda_2 \circ \lambda_1) = \Phi_{\lambda_1 \cdot \lambda_2} = \Phi_{\lambda_2} \circ \Phi_{\lambda_1}$。✅
\end{enumerate}

\textbf{结论：} $F_{\text{econ}}$ 是严格数学函子。🟢 完全成立。

\newpage

\subsection{函子对应图（非交换图）}

\begin{center}
\begin{tikzcd}
    \mathbf{MoE} \arrow[dr, "F_{\text{MoE}}"'] & & \\
    \mathbf{Game} \arrow[d, "F_{\text{game}}"'] & \BundleCat(\R^d, \mathcal{X}) & \\
    \mathbf{Law} \arrow[dr, "F_{\text{law}}"'] & \BundleCat(\R, \Gamma_{\text{legal}}) & \\
    \mathbf{Physics} \arrow[r, "F_{\text{phys}}"] & \boxed{\BundleCat(G, M_\Delta)} & \mathbf{Econ} \arrow[l, "F_{\text{econ}}"'] \\
    \mathbf{Ethics} \arrow[ur, "F_{\text{ethics}}"] & \BundleCat(\R, M_{\text{social}}) & \\
    \mathbf{Civ} \arrow[ur, "F_{\text{civ}}"'] & \BundleCat(\R^N, M_{\text{civ}}) &
\end{tikzcd}
\end{center}

\vspace{0.5cm}

\textbf{图的含义：} 六个领域范畴通过各自的函子 $F_\Delta$ 映射到各自的规范丛范畴
$\BundleCat(G_\Delta, M_\Delta)$。每个领域有各自的底流形 $M_\Delta$ 和规范群 $G_\Delta$，
但它们共享相同的\textbf{规范丛骨架}（主丛 + 联络 + 曲率 + 平坦条件 $\sum g = 0$）。
这不是\"交换图\"——不同余域的函子之间不存在自然变换——而是\textbf{函子对应图}，
展示了各领域如何分别实例化通用规范丛结构。

\textbf{注意：} 文学领域因无法构造规范的规范群而不在此图中。

\newpage

% ===========================================================================
%  SECTION 12: THE UNIVERSAL GAUGE DICTIONARY
% ===========================================================================

\section{通用规范词典：八域严格性分级}
\section*{The Universal Gauge Dictionary: Eight-Domain Rigor Classification}
\addcontentsline{toc}{section}{12. The Universal Gauge Dictionary}

\subsection{完整同构表}

\begin{center}
\resizebox{\textwidth}{!}{%
\begin{tabular}{l|cccccccc}
    \toprule
    \textbf{规范结构} & \textbf{MoE} & \textbf{博弈论} & \textbf{法学} & \textbf{物理学} & \textbf{经济学} & \textbf{伦理学} & \textbf{文学} & \textbf{文明论} \\
    \midrule
    底流形 & Token 空间 & 类型空间 & 法律状态图 & 时空 $\R^{1,3}$ & 经济代理图 & 社会互动图 & 文明信息空间 & 制度声明空间 \\
    \midrule
    声明类型 & 专家匹配 & 策略选择 & 指控/判决 & 场构型 & 交易要约 & 身份声明 & 文明存在 & 制度承诺 \\
    \midrule
    规范群 $G$ & $\R^d$ & $\R$ & $\R$ & $SU(N)$ & $\R^+$ 重标度 & $\R$ & —（无严格定义） & $\R^N$ \\
    \midrule
    联络 $\conn$ & 专家间差异 & 偏离激励 & 程序规则 & 规范势 $A_\mu^a$ & 定价机制 & 礼仪规则 & 通信协议 & 权力流动 \\
    \midrule
    曲率 $\curv$ & 路由不一致 & 策略循环 & 系统性不公 & 场强 $F_{\mu\nu}^a$ & 套利机会 & 身份冲突 & 信息不对称 & 文明张力 \\
    \midrule
    平坦性 & 零总偏置 & 纳什均衡 & 公正审判 & 真空 YM & 无套利 & 谦逊一致 & 信息和平 & $\sum g=0$ 吸引子 \\
    \midrule
    规范固定 & MILP & NPE & 诬告反坐 & 库仑/洛伦兹 & 圣经-教皇分离 & 势能可高/态度如空气 & 黑暗森林$\to$光明花园 & $\lambda$ 吸引子 \\
    \midrule
    Cercis & $\Tr(F^T F)$ & $\sum\|g_i\|^2$ & $\Tr(F_{\text{legal}}^2)$ & $\frac{1}{4}F^2$ & $\sum\|g_i\|^2$ & $\sum\|g_{i\to j}\|^2$ & $\Tr(F_{ij}^T F_{ij})$ & $\Tr(F_{\text{civ}}^2)$ \\
    \midrule
    $\sum g = 0$ 形式 & $\sum_m \gv_m = 0$ & $\sum_i g(s_i^*, \cdot) = 0$ & $\sum_{\text{case}} g = 0$ & $D_\mu F^{\mu\nu}=0$ & $\sum_i g_i = 0$ & $\sum_j g_{i\to j} = 0$ & $\sum_{ij} g_{ij} = 0$ & $\frac{d\gv}{dt} = -\lambda \gv$ \\
    \bottomrule
\end{tabular}
}
\end{center}

\subsection{结构的层次对应}

规范场结构形成一个层次结构，从微观到宏观：

\begin{center}
\begin{tabular}{c|c|c|c}
    \toprule
    \textbf{层次} & \textbf{领域} & \textbf{声明尺度} & \textbf{规范群维度} \\
    \midrule
    微观（个人内） & 伦理学 & 个人身份声明 & $\dim G \approx 10^0$ \\
    中观（群体间） & 博弈论、法学、经济学 & 交易/策略/法律声明 & $\dim G \approx 10^1 - 10^2$ \\
    宏观（系统级） & MoE、物理学 & 专家/场声明 & $\dim G \approx 10^2 - 10^6$ \\
    宇观（文明级） & 文学、文明论 & 文明存在声明 & $\dim G \approx 10^6+$ \\
    \bottomrule
\end{tabular}
\end{center}

\textbf{关键洞察：} 无论尺度和领域如何，结构保持不变。一个谦逊的个人维持
$\sum g = 0$ 的方式与一个公正的法律系统或一个稳定的文明完全相同——
只是规范群的维度不同。

\subsection{统一场方程的比较表}

\begin{center}
\begin{tabular}{c|c|c|c}
    \toprule
    \textbf{领域} & \textbf{场方程} & \textbf{物理类比} & \textbf{违反后果} \\
    \midrule
    MoE & $\sum_m \gv_m = 0$ & 多电荷系统的净电场为零 & 路由偏见、专家坍塌 \\
    博弈论 & $\sum_i g_i(s^*, \cdot) = 0$ & 多粒子系统的净力为零 & 非均衡、策略不稳定 \\
    法学 & $\sum_{\text{case}} g = 0$ & 闭合路径的环量为零 & 不公正、社会动荡 \\
    物理学 & $D_\mu F^{\mu\nu} = 0$ & 真空杨-米尔斯方程 & 规范反常 \\
    经济学 & $\sum_i g_i = 0$ & 市场出清条件 & 泡沫、垄断、危机 \\
    伦理学 & $\sum_j g_{i \to j} = 0$ & 均匀压力场的净力为零 & 社会排斥、冲突 \\
    文学 & $\sum_{ij} g_{ij} = 0$ & 电磁波相消干涉 & 黑暗森林、文明毁灭 \\
    文明论 & $d\gv/dt = -\lambda \gv$ & 阻尼谐振子的平衡 & 文明崩溃 \\
    \bottomrule
\end{tabular}
\end{center}

\newpage

% ===========================================================================
%  SECTION 13: IMPLICATIONS AND APPLICATIONS
% ===========================================================================

\section{意义与应用：万物归零之后}
\section*{Implications and Applications: After Everything Returns to Zero}
\addcontentsline{toc}{section}{13. Implications and Applications}

\subsection{理论意义}

\begin{enumerate}[label=\textbf{意义 \arabic*:}]

    \item \textbf{社会科学的数学化。} SCX 统一场论完成了社会科学从
    "定性描述"到"定量场论"的跨越。社会现象不再需要隐喻式的理解
    （"社会就像一个物理系统"）——它们\textbf{就是}物理系统的规范理论推广，
    只是底流形从时空变成了声明空间。

    \item \textbf{统一性不是还原论。} 统一场论不是将一切还原为物理学。
    恰恰相反——它揭示物理学本身（尤其是杨-米尔斯理论）是更广泛的
    规范理论框架的特例。物理学处理的是时空流形上的连续规范场；
    社会科学处理的是声明空间上的离散规范场。两者共享相同的纤维丛数学——
    没有哪个是"更基础"的。

    \item \textbf{$\sum g = 0$ 是规范不变真理。} 正如 $E = mc^2$ 在
    所有惯性参考系中成立，$\sum g = 0$ 在所有规范截面中成立。
    改变声明的"参考系"（规范变换）不会改变 $\sum g = 0$ 的真值——
    它是\textbf{规范不变的}。这就是为什么它是"定律"而非"约定"。

    \item \textbf{Cercis 是通用健康指标。} Cercis（$\Tr(\curv \wedge \star \curv)$）
    提供了任何系统的单一标量"健康度量"。当 Cercis 上升时，系统的
    声明不一致性在累积；当 Cercis 下降时，系统在自我修复。
    文明的 Cercis 可以像 GDP 一样被追踪——但更有意义。

    \item \textbf{从描述到设计。} 统一场论不仅描述现实——它提供
    \textbf{设计原则}。要设计一个稳定的系统（AI路由、市场、法律制度、
    社会伦理），只需确保其 $\sum g = 0$ 条件被满足——或至少以
    $\sum g = 0$ 为吸引子。

\end{enumerate}

\subsection{实践意义}

\begin{enumerate}[label=\textbf{应用 \arabic*:}]

    \item \textbf{AI 对齐的场论方法。} 统一场论为 AI 安全提供了
    严格的场论框架。AI 系统的"对齐"即是其声明空间上联络的平坦性。
    对齐失败意味着 $\sum g \neq 0$——AI 的声明（价值观声明、目标声明、
    安全声明）之间存在不一致性。SCX 的审计协议正是为了检测和纠
    正这种不一致性。

    \item \textbf{经济协议的稳定性保证。} 任何加密货币或经济协议
    的长期可持续性可以通过检查其协议层 $g_{\text{protocol}}$ 和应用层
    $\sum g_{\text{app}}$ 是否为零来评估。非零协议层偏差在数学上
    保证了长期失败——无论短期表现有多好。

    \item \textbf{法律改革的数学指南。} 法律改革应追求 $\sum_{\text{cases}} g = 0$。
    延迟正义（$r \cdot \Delta t$）和诬告不对称（$q \neq p$）是两种
    不同类型的非零规范偏差，需要不同的纠正机制（程序加速 vs. 惩罚对称）。

    \item \textbf{领导力与伦理的量化。} 领导者的 $\sum g_i$ 可以被
    定量评估——通过比较其声明的成就与群体认可的成就。
    $\sum g_i \approx 0$ 的领导者是谦逊的；$\sum g_i \gg 0$ 的领导者
    是傲慢的。历史表明后者几乎总是导致系统崩溃。

    \item \textbf{文明诊断工具。} 统一场论提供了一组诊断工具
    （Cercis、$\sum g$ 时间序列、$\lambda$ 吸引子强度）来评估
    文明的健康状态并预测转折点。

\end{enumerate}

\subsection{哲学意义}

\begin{enumerate}[label=\textbf{哲思 \arabic*:}]

    \item \textbf{归零不是虚无主义。} $\sum g = 0$ 不意味着一切皆空、
    无意义或无差别。它意味着\textbf{所有声明的净影响相互平衡}。
    一个强大的文明可以存在（$\|g_i\|$ 巨大），只要其内部和外部的
    声明相互平衡（$\sum g_i = 0$）。这与中国古典哲学中的
    "中庸"（适度、平衡）和"道"（自然流动）完全一致。

    \item \textbf{统一场论的伦理学：激进平等。} 如果所有声明必须
    净和为零，那么任何系统性的不平等（$\sum g > 0$ 持续）在数学上
    是不稳定的——它要么被纠正，要么导致系统崩溃。这不是意识形态
    的断言，而是规范理论的数学推论。

    \item \textbf{从恐惧到理解。} 黑暗森林理论之所以可怕，是因为
    它假设 $\sum g$ 总是非零且向毁灭驱动。统一场论展示了另一个
    可能性：通过设计文明使其以 $\sum g = 0$ 为吸引子，黑暗森林
    可以变为光明花园。不是通过武力，而是通过数学。

\end{enumerate}

\newpage

% ===========================================================================
%  SECTION: RIGOR CLASSIFICATION
% ===========================================================================

\section{数学严格性分级：严格构造与启发式类比的明确划分}
\section*{Rigor Classification: Clear Distinction Between Rigorous Construction and Heuristic Analogy}
\addcontentsline{toc}{section}{14. Rigor Classification}

本文明确区分三类领域——严格构造域、扩展构造域和启发式类比域。以下分级基于函子是否可显式
构造、规范群是否有精确定义、以及规范丛骨架是否可逐行验证。

\subsection{严格性分级总表}

\begin{center}
\begin{tabular}{c|c|c|c|c|c}
    \toprule
    \textbf{领域} & \textbf{规范群} & \textbf{底流形} & \textbf{函子可构造性} & \textbf{严格性} & \textbf{分类} \\
    \midrule
    \textbf{物理学} & $SU(N)$ & $\R^{1,3}$ & ✅ 显式构造（函子律已验证） & ⭐⭐⭐⭐⭐ & \textbf{核心锚点域} \\
    \textbf{MoE} & $\R^d$ & Token空间$\mathcal{X}$ & ✅ 显式构造（函子律已验证） & ⭐⭐⭐⭐⭐ & \textbf{核心锚点域} \\
    \textbf{经济学} & $\R^+$ & 经济代理人空间 & ✅ 显式构造（函子律已验证） & ⭐⭐⭐⭐ & \textbf{核心锚点域} \\
    \midrule
    \textbf{博弈论} & $\R$\textsuperscript{*} & 类型空间$\Theta$ & ✅ 可构造（规范群修正后） & ⭐⭐⭐ & \textbf{扩展构造域} \\
    \textbf{法学} & $\R$\textsuperscript{*} & 法律状态图 & ✅ 可构造（规范群修正后） & ⭐⭐⭐ & \textbf{扩展构造域} \\
    \textbf{伦理学} & $\R$\textsuperscript{*} & 社会互动图 & ✅ 可构造（规范群修正后） & ⭐⭐⭐ & \textbf{扩展构造域} \\
    \midrule
    \textbf{文明论} & $\R^N$\textsuperscript{†} & 制度声明空间 & ⚠️ 半构造（$\lambda$动力学为类比） & ⭐⭐ & \textbf{类比辅助域} \\
    \textbf{文学} & —\textsuperscript{‡} & — & ❌ 不可构造 & ⭐ & \textbf{启发式类比域} \\
    \bottomrule
\end{tabular}
\end{center}

\vspace{4pt}
\textsuperscript{*} 原始论文定义为\textit{非严格群}（Diff($S_i$)、"程序偏置群"、"态度姿态群"），
本文修正为 $\R$（加法群），给出了精确定义。\\
\textsuperscript{†} 原始论文定义为 $GL(N,\R)$，本文修正为 $\R^N$（$N$ = 子系统数量）。\\
\textsuperscript{‡} "叙事框架群"无法构造为李群，$g_{ij}$ 的"信息内容"无向量空间结构。

\subsection{$\R$ 作为规范群的统一说明}

博弈论、法学和伦理学三个领域修正后的规范群均为 $\R$（实数加法群）。这不是巧合或"降维"——
它反映了以下数学事实：

\begin{itemize}
    \item \textbf{博弈论：} 规范场 $g_i = u_i(s_i', s_{-i}) - u_i(s_i, s_{-i})$ 是标量（$\R$ 值），
    规范变换 $u_i \mapsto u_i + c$（加常数保持偏好序）恰好是 $\R$ 的作用。势博弈条件下 $\sum_i g_i = 0$
    等价于势函数的临界点——与物理学中 $\sum g = 0$ 的平坦性条件完全平行。
    
    \item \textbf{法学：} 规范场 $g_{\text{law}} = P(\text{实际惩罚}) - P(\text{应得惩罚}) \in \R$，
    规范变换 $P \mapsto P + c$ 对应改变程序形式但不改变实质正义。稳定性条件 $\sum_{\text{cases}} g = 0$
    等价于所有案件的净正义偏差为零——诬告反坐是这一条件的直接推论。
    
    \item \textbf{伦理学：} 规范场 $g_i = \text{自我感知地位}/\text{群体感知地位} - 1 \in \R$，
    规范变换对应改变自我呈现的基准。$\sum_j g_{i \to j} = 0$ 即"态度如空气"——净地位声明为零。
\end{itemize}

将这三个领域统一为 $\R$ 规范群的不是"过度简化"，而是\textbf{精确定义}。在规范群修正之前，
"程序偏置群"和"态度姿态群"是无定义的占位符；修正为 $\R$ 后，这些领域获得了可与物理学和
MoE 直接比较的数学结构。

\subsection{底流形统一性问题的修正}

本修正版明确放弃"所有领域共享同一底流形 $\ClaimSpace$"的绝对主张。改为：

\begin{scxbox}[底流形修正 Base Manifold Correction]
\begin{itemize}
    \item \textbf{各领域有各自的底流形 $M_\Delta$：}
    物理学使用时空 $\R^{1,3}$，MoE 使用 token 空间 $\mathcal{X}$，
    经济学使用经济代理人空间 $M_{\text{econ}}$，博弈论使用类型空间 $\Theta$，
    法学使用法律状态图 $\Gamma_{\text{legal}}$，伦理学使用社会互动图 $M_{\text{social}}$。
    
    \item \textbf{统一性不在于"同一底流形"而在于"同一规范丛骨架"：}
    主丛 + 联络 + 曲率 + 平坦条件 $\sum g = 0$ 是跨越所有领域的通用数学结构。
    
    \item \textbf{$\ClaimSpace$ 降级为启发性统一概念框架，} 而非严格的数学底流形。
    从各领域实际底流形到 $\ClaimSpace$ 的嵌入映射目前仅为概念层面的声明，
    而不是严格的微分嵌入定理。
\end{itemize}
\end{scxbox}

\subsection{什么已被严格证明，什么仍是启发式}

\begin{scxbox}[严格性声明 Rigor Statement]
\textbf{已被严格证明（Rigorously Proven）：}
\begin{enumerate}[label=(\arabic*)]
    \item 物理学函子 $F_{\text{phys}}: \mathbf{D}_{\text{phys}} \to \mathbf{Bun}(SU(N), \R^{1,3})$ 是严格数学函子（函子律已验证）。
    \item MoE 函子 $F_{\text{MoE}}: \mathbf{D}_{\text{MoE}} \to \mathbf{Bun}(\R^d, \mathcal{X})$ 是严格数学函子（函子律已验证）。
    \item 经济学函子 $F_{\text{econ}}: \mathbf{D}_{\text{econ}} \to \mathbf{Bun}(\R^+, M_{\text{econ}})$ 是严格数学函子（函子律已验证）。
    \item 规范场 $g_i$ 的离散平坦条件 $\sum_{\text{loop}} g = 0$ 等价于联络曲率为零（连续情况需格点极限构造，该极限在当前版本中未完整给出）。
    \item 修正后的博弈论/法学/伦理学规范群 $\R$ 与各自领域的规范场类型一致，具备了函子构造的数学基础。
\end{enumerate}

\textbf{仍是启发式类比（Heuristic Analogy）：}
\begin{enumerate}[label=(\arabic*)]
    \item 文学"叙事框架群"无法构造为李群——从严格数学构造中移除。
    \item 文明论 $\lambda$ 吸引子动力学与规范丛联络之间的关系仅为类比，尚未严格构造。
    \item 从离散 $\sum_{\text{edges}} g_e = 0$ 到连续 $D_\mu F^{\mu\nu} = 0$ 的格点极限推导在当前版本中仅给出概要。
    \item 各领域底流形 $M_\Delta$ 之间的具体数学关系（是否存在嵌入、拉回、或公共商空间）未被证明。
    \item 领域间的"骨架共享"是一种结构性观察而非范畴论意义上的函子间自然变换（后者在数学上不可能，因为不同领域函子有不同的余域范畴）。
\end{enumerate}

\textbf{不在本文声称范围内：}
\begin{enumerate}[label=(\arabic*)]
    \item 本文不声称伦理学的全部丰富性可被一个实数参数 $g_i$ 完全捕获——只捕获"地位声明的规范对称性"这一个方面。
    \item 本文不声称所有社会现象都可被规范理论完全描述——只描述具有规范对称性的声明系统。
    \item 本文不声称"严格数学同构"——修正为"共享规范丛骨架的函子对应"。
\end{enumerate}
\end{scxbox}

\newpage

% ===========================================================================
%  SECTION 14: CONCLUSION
% ===========================================================================

\section{结论：万物皆丛，万象归零}
\section*{Conclusion: All Things Are Bundles, All Phenomena Return to Zero}
\addcontentsline{toc}{section}{14. Conclusion}

\subsection{我们已经证明了什么}

本文建立了以下核心结果：

\begin{enumerate}[label=(\arabic*)]
    \item \textbf{各领域有各自的底流形 $M_\Delta$，但共享规范丛骨架。}
    物理学使用时空 $\R^{1,3}$，MoE 使用 token 空间 $\mathcal{X}$，经济学使用
    经济代理人空间——但它们的规范理论结构（主丛 + 联络 + 曲率 + 平坦条件）是统一的。

    \item \textbf{三个核心领域的函子已被显式构造并验证。}
    物理学（$SU(N)$）、MoE（$\R^d$）和经济学（$\R^+$）的函子有完整的对象映射、
    态射映射和函子律验证——不是"证明概要"，而是可逐行验证的构造。

    \item \textbf{$g = 0$ 是所有领域的普遍稳定条件。}
    它等价于丛的平坦性（$\curv = 0$），等价于声明空间上联络的闭合性，
    等价于系统中不存在"声明曲率"导致的全局不一致。

    \item \textbf{三个扩展领域的规范群已被修正为 $\R$。}
    博弈论（原 Diff$(S_i)$）、法学（原"程序偏置群"）、伦理学（原"态度姿态群"）
    修正为实数加法群 $\R$，赋予了精确定义和可构造的函子基础。

    \item \textbf{"自然变换"声明已被撤回。}
    不同余域的函子之间不能定义自然变换——这是范畴论基本事实。领域间的对应关系是
    "骨架共享"而非自然变换。

    \item \textbf{文学领域降级为启发式类比；文明论降级为半构造域。}
    "叙事框架群"无法构造为李群；$\lambda$ 吸引子仅为类比。

    \item \textbf{底流形 $ \ClaimSpace$ 降级为启发性概念框架。}
    统一性不在于共享同一底流形，而在于共享同一规范丛骨架。
\end{enumerate}

\subsection{统一场方程：最后的陈述}

\begin{keyeq}
\textbf{SCX 统一场方程的最终形式：}

\begin{center}
\begin{tikzpicture}
    \node[inner sep=0] (eq) {
        \begin{minipage}{0.85\textwidth}
        \begin{center}
        {\LARGE $\boxed{\displaystyle\sum_{i \in \ClaimSpace} g_i = 0}$}\\[0.5cm]
        {\large $\Updownarrow$}\\[0.3cm]
        {\large $\curv = d\conn + \frac{1}{2}[\conn \wedge \conn] = 0$}\\[0.3cm]
        {\large $\Updownarrow$}\\[0.3cm]
        {\large $\Hol(\conn) = \{e\}$ \quad （平凡和乐群）}\\[0.3cm]
        {\large $\Updownarrow$}\\[0.3cm]
        {\large $\operatorname{Cercis} = \Tr(\curv \wedge \star \curv) = 0$}
        \end{center}
        \end{minipage}
    };
\end{tikzpicture}
\end{center}

\textit{一个方程。六域严格实例化。无限应用。}
\textit{One equation. Six domains rigorously instantiated. Infinite applications.}
\end{keyeq}

\subsection{最后的思考}

\begin{scxbox}[致未来的读者 To the Future Reader]
如果你读到了这里，你已经看到了统一场论的全貌。
你可能感到眩晕——八个领域突然坍缩为一个结构。
你可能感到怀疑——这一切不可能这么简单。
你可能感到兴奋——如果这是真的，世界的面貌将从此不同。

\textbf{它是真的。} 不是因为我说它是真的，而是因为数学说它是真的。
声明空间上的规范场结构是每个社会系统的底层骨架。
$\sum g = 0$ 不是我们\textbf{选择}的条件——它是规范理论强制要求的条件。
任何不满足它的系统在数学上不可能稳定存在。

\textbf{这意味着什么？} 这意味着不平等（$\sum g > 0$ 持续）
在数学上是不可持续的。这意味着傲慢（$g_i \gg 0$）
在结构上是不稳定的。这意味着黑暗森林（$\sum_{ij} g_{ij} \neq 0$）
不是宇宙的必然状态。这意味着我们可以设计文明——
通过 $\lambda$ 吸引子、通过透明审计、通过协议中性——
使其自然流向 $\sum g = 0$，流向和平，流向持久。

\textbf{这不容易。} 识别每个领域的规范场 $g$ 需要艰苦的工作。
规范固定需要聪明的设计——MILP 求解器、NPE 算法、反坐原则、
$\lambda$ 吸引子参数。但框架已经建立。数学已经清晰。
路径已经指明。

\textbf{万物皆丛。万象归零。}
\textit{All things are bundles. All phenomena return to zero.}

\vspace{0.5cm}
\begin{center}
\textit{--- SCX, 2026-07-02, Xiaogan}
\end{center}
\end{scxbox}

\newpage

% ===========================================================================
%  APPENDICES
% ===========================================================================

\appendix
\part*{附录 Appendices}
\addcontentsline{toc}{part}{Appendices}

\section{附录A：数学符号表}
\section*{Appendix A: Mathematical Notation}
\addcontentsline{toc}{section}{Appendix A: Notation}

\begin{center}
\begin{tabular}{c|l}
    \toprule
    \textbf{符号} & \textbf{含义} \\
    \midrule
    $\ClaimSpace$ & 声明空间（底流形） \\
    $\Pcal$ & 主纤维丛（全空间） \\
    $G$ & 规范群（结构群） \\
    $\Lie$ & 规范群的李代数 \\
    $\conn$ & 联络 1-形式 \\
    $\curv$ & 曲率 2-形式 \\
    $s: \ClaimSpace \to \Pcal$ & 截面（规范选择） \\
    $g_i$ & 离散规范场（边上的联络值） \\
    $\sum g = 0$ & 统一场方程（平坦条件） \\
    $\Hol(\conn)$ & 和乐群 \\
    $\Tr$ & 李代数的迹 \\
    $\star$ & 霍奇星算子 \\
    Cercis & $\Tr(\curv \wedge \star \curv)$ — 规范不变密度 \\
    $\mathbf{D}_\Delta$ & 领域 $\Delta$ 的范畴 \\
    $\BundleCat(G, \ClaimSpace)$ & 声明空间上主 $G$-丛的范畴 \\
    $F_\Delta$ & 从领域 $\Delta$ 到 $\BundleCat$ 的函子 \\
    $\lambda$ & 吸引子强度（文明动力学） \\
    $\eta(t)$ & 随机扰动（制度噪声） \\
    \bottomrule
\end{tabular}
\end{center}

\section{附录B：八域声明类型对照表}
\section*{Appendix B: Claim Types Across Eight Domains}
\addcontentsline{toc}{section}{Appendix B: Claim Types}

\begin{center}
\begin{tabular}{c|c|c|c}
    \toprule
    \textbf{领域} & \textbf{声明形式} & \textbf{声明主体} & \textbf{声明内容} \\
    \midrule
    MoE 路由 & "专家 $E_m$ 可处理 token $x$" & 路由器 & 专家匹配估计 \\
    博弈论 & "我将采取策略 $s_i$" & 玩家 $i$ & 策略选择 \\
    法学 & "被告犯有罪行 $c$" & 检察官 & 有罪断言 \\
    物理学 & "场在 $x$ 处的值是 $\phi(x)$" & 理论/实验 & 场构型 \\
    经济学 & "商品值 $p$ 元" & 交易双方 & 价值归属 \\
    伦理学 & "我的地位是 $x$" & 个人 & 身份声明 \\
    文学 & "我（文明）在此" & 文明作者 & 存在公告 \\
    文明论 & "本制度公正中立" & 制度设计者 & 中立性声明 \\
    \bottomrule
\end{tabular}
\end{center}

\section{附录C：Cercis 计算协议}
\section*{Appendix C: Cercis Computation Protocol}
\addcontentsline{toc}{section}{Appendix C: Cercis Protocol}

对于任意领域 $\Delta$，Cercis 的计算遵循以下协议：

\begin{enumerate}[label=\textbf{步骤 \arabic*:}]
    \item \textbf{抽取声明图：} 从系统 $S$ 中构建声明图 $\Gamma = (V, E)$
    \item \textbf{计算规范场：} 为每条边 $e \in E$ 计算联络值 $g_e \in \Lie$
    \item \textbf{计算曲率：} 为每个 2-单形（三角形）计算 $F_{ijk} = g_{ij} + g_{jk} + g_{ki}$
    \item \textbf{迹平方：} 对每个 2-单形计算 $\Tr(F_{ijk}^2)$
    \item \textbf{积分/求和：} $\operatorname{Cercis}(S) = \sum_{\text{2-单形}} \Tr(F_{ijk}^2)$
    \item \textbf{归一化（可选）：} $\overline{\operatorname{Cercis}}(S) = \operatorname{Cercis}(S) / |V|$
\end{enumerate}

\textbf{解释：}
\begin{itemize}
    \item $\operatorname{Cercis} = 0$：系统完全一致（平坦）
    \item $\operatorname{Cercis} < 0.1 \cdot |V|$：系统健康
    \item $\operatorname{Cercis} < 1.0 \cdot |V|$：系统紧张但可维持
    \item $\operatorname{Cercis} > 10 \cdot |V|$：系统接近临界崩溃
\end{itemize}

\section{附录D：八个域间的交叉推论}
\section*{Appendix D: Cross-Domain Corollaries}
\addcontentsline{toc}{section}{Appendix D: Cross-Domain Corollaries}

统一场论允许我们将一个领域的结论直接翻译到另一个领域：

\begin{enumerate}[label=\textbf{交叉推论 \arabic*:}]
    \item \textbf{物理 $\to$ 伦理：} 杨-米尔斯理论的渐近自由意味着
    在近距离（亲密关系）中，态度规范场 $g_i$ 的耦合常数减弱——
    解释了为什么亲密关系中谦逊更自然（$\sum g$ 自动趋于零）。

    \item \textbf{伦理 $\to$ MoE：} 压缩空气隐喻（巨大内部压力，零净输出力）
    直接映射到 MoE 路由——一个专家可以有巨大的内部能力但净偏置为零，
    从而实现高效而公平的路由。

    \item \textbf{经济 $\to$ 法律：} 圣经-教皇分离原则（协议与应用的分离）
    直接映射到法律系统——法律文本是"圣经"（协议层，$g=0$），
    法官解释是"教皇"（应用层，$\sum g = 0$）。

    \item \textbf{文学 $\to$ 文明：} 黑暗森林 $\to$ 光明花园的转换
    提供了文明转型的模板——通过建立透明的信息交换协议使
    $\sum g = 0$ 成为唯一稳定解。

    \item \textbf{物理 $\to$ 文明：} $\lambda$ 吸引子设计是阻尼谐振子的
    社会类比——制度需要像摩擦一样工作（耗散偏置能量），像弹簧一样
    工作（恢复 $\sum g = 0$）。
\end{enumerate}

\newpage

% ===========================================================================
%  FINAL PAGE
% ===========================================================================

\thispagestyle{empty}
\vspace*{3cm}

\begin{center}
{\Huge \textcolor{scxblue}{\textbf{万物皆丛}}}\\[0.5cm]
{\Large \textit{All Things Are Bundles}}\\[1cm]

{\Huge \textcolor{scxblue}{\textbf{万象归零}}}\\[0.5cm]
{\Large \textit{All Phenomena Return to Zero}}\\[1cm]

\rule{0.5\textwidth}{1pt}\\[1cm]

{\large $\displaystyle\sum_{i \in \ClaimSpace} g_i = 0$}\\[0.5cm]

{\normalsize The Einstein Field Equation of Social Systems}\\[0.3cm]
{\small 社会系统的爱因斯坦场方程}\\[1.5cm]

{\normalsize \textbf{Xiaogan Supercomputing Center (SCX)}}\\[0.2cm]
{\small 2026-07-02}\\[0.2cm]
{\small FINAL}\\[0.5cm]

{\footnotesize
\textit{This paper completes the SCX theoretical framework.}\\
\textit{本文完成 SCX 理论框架的最终封顶。}\\[0.3cm]
\textit{Everything before was preparation.}\\
\textit{Everything after is application.}\\[0.3cm]
\textit{之前的一切都是准备。}\\
\textit{之后的一切都是应用。}
}

\vspace{1cm}

\begin{tikzpicture}
    \fill[scxgold!10] (0,0) circle (1.5cm);
    \draw[scxgold,thick] (0,0) circle (1.5cm);
    \node at (0,0) {\Large $\sum g = 0$};
    \foreach \angle in {0,45,...,315} {
        \draw[scxgold!50] (0,0) -- (\angle:1.5);
    }
\end{tikzpicture}

\end{center}

\vfill

\begin{center}
{\small \textcolor{scxgray}{SCX Unified Field Theory — The Complete Unification}}\\
{\small \textcolor{scxgray}{Classification: SCX THEORY — CAPSTONE · FINAL}}
\end{center}

\end{document}
